% concatenated from enhanced.tex
\documentclass[12pt]{amsart}
\usepackage{tikz-cd} 
\usepackage{xy}
\usepackage{graphicx}
\usepackage{subfigure}
\usepackage{amsmath,amssymb,amsthm,amsfonts,amsrefs}
\usepackage{newclude}
\usepackage{multirow}
\usepackage{dynkin-diagrams}
\counterwithin{equation}{section}
\newtheorem{lm}{Lemma}[section]
\newtheorem{theorem}[lm]{Theorem}
\newtheorem{cor}[lm]{Corollary}
\newtheorem{prop}[lm]{Proposition}
\newtheorem{conj}{Conjecture}[section]
\theoremstyle{definition}
\newtheorem{defn}[lm]{Definition}
\newtheorem{exa}[lm]{Example}
\newtheorem{exe}{Exercise}[section]
\theoremstyle{remark}\newtheorem{rmk}[lm]{Remark}
\usepackage{latexsym}
\usepackage{CJKutf8}
\usepackage{indentfirst}
\usepackage{mathrsfs}
\usepackage{geometry}
%\geometry{inner=1in,outer=1.25in,vmargin=1in}
\usepackage{imakeidx}
\usepackage{appendix}
\usepackage[colorlinks=true,hyperfootnotes=true]{hyperref}
\usepackage{mathdots}
\usepackage{arydshln}

\newcommand{\BA}{{\mathbb {A}}}
\newcommand{\BB}{{\mathbb {B}}}
\newcommand{\BC}{{\mathbb {C}}}
\newcommand{\BD}{{\mathbb {D}}}
\newcommand{\BE}{{\mathbb {E}}}
\newcommand{\BF}{{\mathbb {F}}}
\newcommand{\BG}{{\mathbb {G}}}
\newcommand{\BH}{{\mathbb {H}}}
\newcommand{\BI}{{\mathbb {I}}}
\newcommand{\BJ}{{\mathbb {J}}}
\newcommand{\BK}{{\mathbb {K}}}
\newcommand{\BL}{{\mathbb {L}}}
\newcommand{\BM}{{\mathbb {M}}}
\newcommand{\BN}{{\mathbb {N}}}
\newcommand{\BO}{{\mathbb {O}}}
\newcommand{\BP}{{\mathbb {P}}}
\newcommand{\BQ}{{\mathbb {Q}}}
\newcommand{\BR}{{\mathbb {R}}}
\newcommand{\BS}{{\mathbb {S}}}
\newcommand{\BT}{{\mathbb {T}}}
\newcommand{\BU}{{\mathbb {U}}}
\newcommand{\BV}{{\mathbb {V}}}
\newcommand{\BW}{{\mathbb {W}}}
\newcommand{\BX}{{\mathbb {X}}}
\newcommand{\BY}{{\mathbb {Y}}}
\newcommand{\BZ}{{\mathbb {Z}}}
\newcommand{\Bk}{{\mathbf {k}}}

\newcommand{\CA}{{\mathcal {A}}}
\newcommand{\CB}{{\mathcal {B}}}
\newcommand{\CC}{{\mathcal {C}}}
\newcommand{\CE}{{\mathcal {E}}}
\newcommand{\CF}{{\mathcal {F}}}
\newcommand{\CG}{{\mathcal {G}}}
\newcommand{\CH}{{\mathcal {H}}}
\newcommand{\CI}{{\mathcal {I}}}
\newcommand{\CJ}{{\mathcal {J}}}
\newcommand{\CK}{{\mathcal {K}}}
\newcommand{\CL}{{\mathcal {L}}}
\newcommand{\CM}{{\mathcal {M}}}
\newcommand{\CN}{{\mathcal {N}}}
\newcommand{\CO}{{\mathcal {O}}}
\newcommand{\CP}{{\mathcal {P}}}
\newcommand{\CQ}{{\mathcal {Q}}}
\newcommand{\CR}{{\mathcal {R}}}
\newcommand{\CS}{{\mathcal {S}}}
\newcommand{\CT}{{\mathcal {T}}}
\newcommand{\CU}{{\mathcal {U}}}
\newcommand{\CV}{{\mathcal {V}}}
\newcommand{\CW}{{\mathcal {W}}}
\newcommand{\CX}{{\mathcal {X}}}
\newcommand{\CY}{{\mathcal {Y}}}
\newcommand{\CZ}{{\mathcal {Z}}}

\newcommand{\RA}{{\mathrm {A}}}
\newcommand{\RB}{{\mathrm {B}}}
\newcommand{\RC}{{\mathrm {C}}}
\newcommand{\RD}{{\mathrm {D}}}
\newcommand{\RE}{{\mathrm {E}}}
\newcommand{\RF}{{\mathrm {F}}}
\newcommand{\RG}{{\mathrm {G}}}
\newcommand{\RH}{{\mathrm {H}}}
\newcommand{\RI}{{\mathrm {I}}}
\newcommand{\RJ}{{\mathrm {J}}}
\newcommand{\RK}{{\mathrm {K}}}
\newcommand{\RL}{{\mathrm {L}}}
\newcommand{\RM}{{\mathrm {M}}}
\newcommand{\RN}{{\mathrm {N}}}
\newcommand{\RO}{{\mathrm {O}}}
\newcommand{\RP}{{\mathrm {P}}}
\newcommand{\RQ}{{\mathrm {Q}}}
\newcommand{\RR}{{\mathrm {R}}}
\newcommand{\RS}{{\mathrm {S}}}
\newcommand{\RT}{{\mathrm {T}}}
\newcommand{\RU}{{\mathrm {U}}}
\newcommand{\RV}{{\mathrm {V}}}
\newcommand{\RW}{{\mathrm {W}}}
\newcommand{\RX}{{\mathrm {X}}}
\newcommand{\RY}{{\mathrm {Y}}}
\newcommand{\RZ}{{\mathrm {Z}}}

\newcommand{\Ad}{{\mathrm{Ad}}}
\newcommand{\Aut}{{\mathrm{Aut}}}

\newcommand{\disc}{{\mathrm{disc}}}
\newcommand{\Div}{{\mathrm{Div}}}
\renewcommand{\div}{{\mathrm{div}}}
\newcommand{\End}{{\mathrm{End}}}

\newcommand{\Frob}{{\mathrm{Frob}}}


\newcommand{\Gal}{{\mathrm{Gal}}}
\newcommand{\GL}{{\mathrm{GL}}}
\newcommand{\Gr}{{\mathrm{Gr}}}
\newcommand{\GSp}{{\mathrm{GSp}}}

\newcommand{\Hom}{{\mathrm{Hom}}}

\renewcommand{\Im}{{\mathrm{Im}}}
\newcommand{\Ind}{{\mathrm{Ind}}}
\newcommand{\inv}{{\mathrm{inv}}}
\newcommand{\Isom}{{\mathrm{Isom}}}
\newcommand{\id}{{\mathrm{id}}}
\newcommand{\Int}{\mathrm{Int}}
\newcommand{\Inn}{\mathrm{Inn}}
\newcommand{\ind}{\mathrm{ind}}

\newcommand{\Jac}{{\mathrm{Jac}}}

\newcommand{\Lie}{{\mathrm{Lie}}}
\newcommand{\Mat}{{\mathrm{Mat}}}
\newcommand{\PM}{{\mathrm{PM}}}
\newcommand{\Mp}{{\mathrm{Mp}}}

\newcommand{\ord}{{\mathrm{ord}}}

\newcommand{\pr}{{\mathrm{pr}}}
\newcommand{\rank}{{\mathrm{rank}}}


\newcommand{\PGL}{{\mathrm{PGL}}}
\newcommand{\Pic}{\mathrm{Pic}}

\renewcommand{\Re}{{\mathrm{Re}}}
\newcommand{\reg}{{\mathrm{reg}}}
\newcommand{\Res}{{\mathrm{Res}}}
\newcommand{\res}{{\mathrm{res}}}
\newcommand{\re}{{\mathrm{e}}}
\newcommand{\rf}{{\mathrm{f}}}
\newcommand{\rk}{{\mathrm{k}}}
\newcommand{\rn}{\mathrm{n}}
\newcommand{\rr}{{\mathrm{r}}}
\newcommand{\rs}{\mathrm{s}}
\newcommand{\rh}{{\mathrm{h}}}


\newcommand{\Sel}{{\mathrm{Sel}}}
\newcommand{\Sim}{{\mathrm{Sim}}}
\newcommand{\SL}{{\mathrm{SL}}}
\newcommand{\PSL}{{\mathrm{PSL}}}
\newcommand{\PGSp}{{\mathrm{PGSp}}}
\newcommand{\Spec}{{\mathrm{Spec}}}
\newcommand{\SO}{{\mathrm{SO}}}
\newcommand{\GO}{{\mathrm{GO}}}
\newcommand{\GSO}{{\mathrm{GSO}}}
\newcommand{\SU}{{\mathrm{SU}}}
\newcommand{\GU}{{\mathrm{GU}}}
\newcommand{\Sym}{{\mathrm{Sym}}}
\newcommand{\symm}{\mathrm{symm}}
\newcommand{\Symb}{{\mathrm{Symb}}}
\newcommand{\BSymb}{{\mathrm{BSymb}}}
%\newcommand{\sgn}{{\mathrm{sgn}}}
\newcommand{\Sp}{{\mathrm{Sp}}}
\newcommand{\Stab}{{\mathrm{Stab}}}
\newcommand{\tr}{{\mathrm{tr}}}
\newcommand{\DO}{{\mathrm{DO}}}
\newcommand{\ur}{{\mathrm{ur}}}
\newcommand{\vol}{{\mathrm{vol}}}
\newcommand{\ab}{{\mathrm{ab}}}
\newcommand{\alg}{\mathrm{alg}}
\newcommand{\Comm}{\mathrm{Comm}}
\newcommand{\Par}{\mathrm{Par}}
\newcommand{\Exp}{\mathrm{Exp}}
\newcommand{\ct}{\mathrm{ct}}
\newcommand{\sm}{\mathrm{sm}}
\newcommand{\triv}{\mathrm{triv}}
\newcommand{\rd}{\mathrm{d}}

\newcommand{\gr}{\operatorname{gr}}
\newcommand{\ad}{\operatorname{ad}}
\newcommand{\dif}{\operatorname{d}\!}
\newcommand{\pro}{\operatorname{pro}}
\newcommand{\Tr}{\operatorname{Tr}}
\newcommand{\Norm}{\operatorname{Norm}}
\newcommand{\Ext}{\operatorname{Ext}}
\newcommand{\otop}{\operatorname{top}}
\newcommand{\diag}{\operatorname{diag}}
\newcommand{\sgn}{\operatorname{sgn}}
\newcommand{\im}{\operatorname{im}}
\newcommand{\Span}{\operatorname{Span}}

\newcommand{\g}{\mathfrak g}
\newcommand{\gC}{{\mathfrak g}_{\C}}
\renewcommand{\k}{\mathfrak k}
\newcommand{\h}{\mathfrak h}
\newcommand{\fH}{\mathfrak H}
\newcommand{\p}{\mathfrak p}
\newcommand{\fq}{\mathfrak q}
\renewcommand{\a}{\mathfrak a}
\renewcommand{\b}{\mathfrak b}
\renewcommand{\c}{\mathfrak c}
\newcommand{\n}{\mathfrak n}
\renewcommand{\u}{\mathfrak u}
%\renewcommand{\v}{\mathfrak v}
\newcommand{\e}{\mathfrak e}
\newcommand{\f}{\mathfrak f}
\renewcommand{\l}{\mathfrak l}
\renewcommand{\t}{\mathfrak t}
\newcommand{\s}{\mathfrak s}
\renewcommand{\r}{\mathfrak r}
\renewcommand{\o}{\mathfrak o}
\newcommand{\m}{\mathfrak m}
\newcommand{\z}{\mathfrak z}
\newcommand{\su}{\mathfrak s \mathfrak u}
\renewcommand{\sl}{\mathfrak s \mathfrak l}
\newcommand{\gl}{\mathfrak g \mathfrak l}
\renewcommand{\sp}{\mathfrak s \mathfrak p}
\newcommand{\fM}{\mathfrak M}
\newcommand{\fN}{\mathfrak N}
\newcommand{\fZ}{\mathfrak Z}
\newcommand{\fX}{\mathfrak X}
\newcommand{\fG}{\mathfrak G}
\newcommand{\fF}{\mathfrak F}
\newcommand{\fS}{\mathfrak S}
\newcommand{\fD}{\mathfrak D}
\newcommand{\fR}{\mathfrak R}
\newcommand{\fI}{\mathfrak I}

\renewcommand{\rk}{\mathrm k}

\newcommand{\Z}{\mathbb{Z}}
\newcommand{\C}{\mathbb{F}}
\newcommand{\R}{\mathbb R}
\renewcommand{\H}{\mathbb{H}}
%\newcommand{\N}{\mathbb{N}}
\newcommand{\K}{\mathbb{K}}
%\renewcommand{\S}{\mathbf S}
\newcommand{\M}{\mathbf{M}}
\newcommand{\A}{\mathbf{A}}
\newcommand{\B}{\mathbf{B}}
\newcommand{\G}{\mathbf{G}}
\newcommand{\V}{\mathbf{V}}
\newcommand{\W}{\mathbf{W}}
\newcommand{\F}{\mathbf{F}}
\newcommand{\E}{\mathbf{E}}
\newcommand{\II}{\mathbf{I}}
\newcommand{\J}{\mathbf{J}}
\renewcommand{\H}{\mathbf{H}}
\newcommand{\X}{\mathbf{X}}
\newcommand{\Y}{\mathbf{Y}}
%\newcommand{\RR}{\mathcal R}
%\newcommand{\FF}{\mathcal F}
%\newcommand{\BB}{\mathcal B}
\newcommand{\HH}{\mathscr H}
\newcommand{\bH}{\mathbb H}
\newcommand{\UU}{\mathcal U}
\newcommand{\MM}{\mathcal M}
%\newcommand{\CC}{\mathcal C}
\newcommand{\DD}{\mathcal D}
\newcommand{\ZZ}{\mathcal Z}
\newcommand{\ve}{{\vee}}
\newcommand{\aut}{\mathcal A}
\newcommand{\abs}[1]{\lvert#1\rvert}
\newcommand{\be}{\mathbf{e}}
\newcommand{\ii}{\mathbf{i}}
\newcommand{\jj}{\mathbf{j}}
\newcommand{\kk}{\mathbf{k}}
\newcommand{\bp}{\mathbf{p}}
\newcommand{\bq}{\mathbf{q}}
\newcommand{\bl}{\mathbf{l}}
\newcommand{\bm}{\mathbf{m}}
\newcommand{\bn}{\mathbf{n}}
\newcommand{\bt}{\mathbf{t}}

\newcommand{\dR}{\mathrm{dR}}
\newcommand{\supp}{\mathrm{supp}}
\newcommand{\Spe}{\mathrm{Sp\'e}}
\newcommand{\la}{\langle\!\langle}
\newcommand{\ra}{\rangle\!\rangle}
\newcommand{\bac}{\backslash\!\!\backslash}
\renewcommand{\mid}{\,:\,}
\newcommand{\tran}[1]{{}^{\mathtt{t}}\!{#1}}
\newcommand{\sW}{\mathsf{W}}
\newcommand{\sD}{\mathscr D}
\newcommand{\tL}{\mathtt L}
\newcommand{\tT}{\mathtt T}
\newcommand{\sM}{\mathscr M}

\newcommand{\cf}{\emph{cf.}~}

\setcounter{tocdepth}{1}

\title[Closed orbits and descents]{Closed Orbits and Descents for Enhanced Standard Representations of Classical Groups}
\date{}
\author{Chen Liang}
\address{School of Mathematical Sciences, Zhejiang University\\
	Hangzhou, 310058, China}\email{12235009@zju.edu.cn}

\begin{document}
	\maketitle
	\begin{abstract} 
    Let $G=\GL_n(\C)$, $\RO_n(\C)$, or $\Sp_{2n}(\C)$ be one of the classical groups over an algebraically closed field $\C$ of characteristic $0$, let $\breve{G}$ be  the MVW-extension  of $G$,  and let $\g$ be the Lie algebra of $G$. 
    In this paper, we classify the closed orbits in the enhanced standard representation $\g\times E$ of $G$, 
    where $E$ is the natural representation if $G=\RO_n(\C)$ or $\Sp_{2n}(\C)$, and is the direct sum of the natural representation and its dual if $G=\GL_n(\C)$.
 Additionally, for every closed $G$-orbit in $\g\times E$, we prove that it is $\breve{G}$-stable, and determine explicitly the corresponding stabilizer group as well as the action on the normal space. 
	\end{abstract}
	\tableofcontents
	\section{Introduction and main results}
    Throughout this paper, let $\C$ be an algebraically closed field of characteristic 0.
    \subsection{Motivations}
Let $\G$ be a reductive algebraic group over a local field $F$ of characteristic $0$.  Function spaces for rational representations of $\G$ play an important role in the study of representations of $\G(F)$. One classical problem is the ``multiplicity one problem'': Given a larger group $\G'(F)$ and an irreducible smooth admissible representation $\pi$ of $\G'(F)$, is the dimension of $\Hom_{\G(F)}(\pi,\BC)$ at most $1$? By the Gelfand-Kazhdan criterion \cite{GK,SZ1} and the Harish-Chandra descent \cite{AG}, the multiplicity one problem can be reduced to proving that certain $\breve{\G}(F)$-equivariant generalized functions on a rational representation of $\G(F)$ must be $0$, where $\breve{\G}(F)$ is an extension of $\G(F)$ by $\{\pm 1\}$. Another problem is the existence of smooth transfer. Let $V$ be a rational representation of $\G$. Given another reductive group $\G'$ over $F$ and a rational representation $V'$ of $\G'$, suppose that there exists a matching of regular semismiple orbits between $V$ and $V'$. The smooth transfer of a Schwartz function $f$ on $V(F)$ is a Schwartz function $f'$ on $V'(F)$ such that $O_{\gamma}(f)=\Delta(\gamma,\gamma')O_{\gamma'}(f')$, whenever a regular semisimple $\gamma\in V(F)$ matches a regular semisimple $\gamma'\in V'(F)$, where $O_{\gamma}(f),O_{\gamma'}(f')$ are suitably defined orbit integrals, and $\Delta(\gamma,\gamma')$ is the transfer factor.

To study the function spaces of the rational representation $V(F)$, one needs to study the geometry of the action of $\G$ on $V$ first, in particular, the classification of closed orbits, the corresponding stabilizer group and the descendants (see Definition \ref{defn:descent}). In this article, we investigate these geometric properties for the enhanced standard representation of classical groups as well as their MVW-extensions (see \S\,\ref{sec:enrepn} for the definitions). Our results may be applied in the proof of multiplicity one theorem and the existence of smooth transfer. In the proof of multiplicity one theorem \cite{AGb,AGRS,SZ}, as pointed out in \cite{AGb}, by applying \cite{AG}*{Theorem 3.2.1}, one of our results (Theorem \ref{thm:mainmvwsimple}) can provide a direct proof of \cite{AGb}*{Proposition 3.2.1}, \cite{AGRS}*{Propositions 3.2, 5.2} and \cite{SZ}*{Propositions 7.1 and 7.2}. In the proof of the existence of smooth transfer \cite{Zha14, Xue19, CZ21} for Jacquet-Rallis relative trace formulas, the action on the general linear side is reduced to the enhanced standard representation of $\GL_n$, and its descendants are used for further reduction.
    
%In this article, we study  geometric properties for the enhanced standard representation of classical groups over $\C$ as well as their MVW-extensions (see \S\,\ref{sec:enrepn} for the definitions).   Our motivation comes from the study of automorphic representations, in particular the relative Langlands program proposed by Ben-Zvi, Sakellaridis and Venkatesh \cite{BZSV}.

%The first motivation  is the problem of smooth transfer. One of aims of the relative Langlands program is to relate certain period integrals with L-functions. Let $(\G,\H_1,\H_2)$ be a triples consisting of a reductive group and two spherical subgroups over a number field $K$. A key method to study period integrals is to compare relative trace formulas for two such triples. The comparison of two relative trace formulas leads to a local question: the existence of smooth transfer. To this end, one needs to study the geometry of the action of $\H_1\times\H_2$ on $\G$, in particular, the classification of regular semisimple orbits. For example, Jacquet–Rallis proposed an approach in \cite{JR} to the global Gan-Gross-Prasad conjecture \cite{GGP} for unitary groups based on comparing the relative trace formulae of two special triples. In the proof of the existence of smooth transfer \cite{Zha14, Xue19, CZ21} for Jacquet-Rallis trace formulas, the action of $\H_1\times\H_2$ on $\G$ in the general linear side is reduced to the enhanced standard representation of $\GL_n$, and its descendants are used to reduce the proof to the lower-rank cases.
    
%	Another motivation comes from  the multiplicty one conjecture. Let $(\G,\H)$ be a pair consisting of a reductive group and a spherical subgroup over a local field $F$ of characteristic $0$. The multiplicity one conjecture asserts that the dimension of $\Hom_{\H(F)}(\pi,\mathbb C)$ is no more than one for every irreducible smooth admissible representation $\pi$ of $\G(F)$, where $\G(F)$ (resp. $\H(F)$) is the group of $F$-rational points of $\G$ (resp. $\H$). By the Gelfand-Kazhdan criterion \cite{GK,SZ1}, the multiplicity one conjecture can be reduced to proving that certain $\breve\H(F)$-equivariant generalized functions on $\g(F)/\h(F)$ must be zero, where $\breve{\H}(F)$ is an extension of $\H(F)$ by $\{\pm 1\}$ and $\g$ (resp. $\h$) is the Lie algebra of $\G$ (resp. $\H$). This leads to the study of the geometry of the action of $\breve\H(F)$ on $\g(F)/\h(F)$, especially the closed orbits and the corresponding stabilizer groups as well as the actions on the normal spaces. For example, in \cite{AGRS, AGb, SZ}, the authors consider the case that $\H$ is a classical group, $\g(F)/\h(F)$ is the sum of the enhanced standard representation and the trivial representation of $\H(F)$, and $\breve{\H}(F)$ is the MVW-extension of $\H(F)$. As pointed out in \cite{AGb}, by applying \cite{AG}*{Theorem 3.2.1}, one of our results (Theorem \ref{thm:mainmvwsimple}) can provide a direct proof of \cite{AGb}*{Proposition 3.2.1}, \cite{AGRS}*{Propositions 3.2, 5.2} and \cite{SZ}*{Propositions 7.1 and 7.2}.

    
    
   
	\subsection{Enhanced standard representations} \label{sec:enrepn}
 Let $G$ be one of the following classical groups:
	\begin{equation}
    \label{eq:classicalgroups}
		\GL_n(\C),\quad\RO_n(\C)\quad\text{and}\quad\Sp_{2n}(\C)\qquad (n\ge 0),
	\end{equation}
	where $\GL_n(\C)$ is the general linear group of rank $n$ over $\C$, 
	\begin{equation*}
		\RO_n(\C)=\left\{g\in\GL_n(\C)\mid g^t\alpha_ng=\alpha_n\right\}
	\end{equation*}
    is the orthogonal group, 
	and
	\begin{equation*}
		\Sp_{2n}(\C)=\left\{g\in\GL_{2n}(\C)\mid g^t\beta_{2n}g=\beta_{2n}\right\}
	\end{equation*}
    is the symplectic group. Here,
    \begin{equation*}
        \alpha_{n}=\left(\begin{matrix}
				&I_{\frac{n}{2}}\\I_{\frac{n}{2}}&
			\end{matrix}\right),\text{ if n even},\quad\alpha_n=
            \left(\begin{matrix}
				1&&\\&&I_{\frac{n-1}{2}}\\&I_{\frac{n-1}{2}}&
			\end{matrix}\right),\text{ if n is odd},
    \end{equation*}
    and
    \begin{equation*}
        \beta_{2n}=\left(\begin{matrix}
			&I_n\\-I_n&
		\end{matrix}\right).
    \end{equation*}
    
    Write $\breve{G}=G\rtimes\{\pm 1\}$ for the \emph{MVW-extension} of $G$ (cf.\,\cite{MVW},\cite{Sun}), where 
    \begin{equation*}
   (-1).g   = \begin{cases*} 
   g^{-t}, & if $G=\GL_n(\C)$,\\
   \alpha_{n}g\alpha_{n},& if $G=\RO_n(\C)$,\\
\alpha_{2n}g\alpha_{2n}, & if $G=\Sp_{2n}(\C)$. 
        \end{cases*}
    \end{equation*}
    In this paper, we view $\breve{G}$ and $G$ as algebraic groups over $\C$. 
            %We often identify $G$ as a subgroup of $\breve{G}$, and regard a rational representation of $\breve{G}$ as a rational representation of $G$ by taking restriction. 

            
	Let $\g=\mathfrak{gl}_n(\C)$, $\mathfrak{o}_n(\C)$, and $\mathfrak{sp}_{2n}(\C)$  be the Lie algebras of $G=\GL_n(\C)$, $\RO_n(\C)$, and $\Sp_{2n}(\C)$, respectively. Define an action of $\breve{G}$ on $\g$ by letting 
    \begin{equation*}
    (g,1).X=gXg^{-1}\quad \text{and}\quad    (g,-1).X=\begin{cases*}
            gX^tg^{-1},&if $G=\GL_n(\C)$,\\
            -g\alpha_{n}X\alpha_{n}g^{-1},&if $G=\RO_n(\C)$,\\
            -g\alpha_{2n}X\alpha_{2n}g^{-1},&if $G=\Sp_{2n}(\C)$.
        \end{cases*}
    \end{equation*}
   Additionally, we define an action of $\breve{G}$ on  the space 
	\begin{equation*}
		E:=\begin{cases*}
			\C^{n\times 1}\times \C^{1\times n},&if $G=\GL_n(\C)$,\\
			\C^{n\times 1},&if $G=\RO_n(\C)$,\\
            \C^{2n\times 1}, &if $G=\Sp_{2n}(\C)$
		\end{cases*}
	\end{equation*}
as follows: 

\begin{equation*}
    (g,\delta).(u,v)=\begin{cases*}
        (gu,vg^{-1}),&if $G=\GL_n(\C)$ and $\delta=1$,\\
        (-gv^t,-u^tg^{-1}),&if $G=\GL_n(\C)$ and $\delta=-1$,
    \end{cases*}
\end{equation*}
and
\begin{equation*}
    (g,\delta).u=\begin{cases*}
        gu,&if $G=\Sp_{2n}(\C)$ or $\RO_n(\C)$ and $\delta=1$,\\
        -g\alpha_{2n}u,&if $G=\Sp_{2n}(\C)$ and $\delta=-1$,\\
        -g\alpha_nu,&if $G=\RO_n(\C)$ and $\delta=-1$
    \end{cases*}
\end{equation*}
    Here, $\C^{p\times q}$ ($p,q\ge 1$) denotes the space of $p\times q$-matrices over $\C$. 
     
	Let $\breve{G}$ act on $\g\times E$ diagonally, so that it is a rational representation of $\breve{G}$ (and $G$ by taking restriction). 
    We call this representation of $\breve{G}$ (resp. $G$) the \emph{enhanced standard representation} of $\breve{G}$ (resp. $G$). 
 The main goal of this paper is to classify the closed $G$-orbits and $\breve{G}$-orbits in  $\g\times E$, and determine their corresponding stabilizer subgroups.

  \subsection{Related works}
     Denote by $\CN_{\g}$ the null cone of $\g$, which consists of all nilpotent matrices in $\g$. The closed $G$-orbits in $\g$ as well as the $G$-orbits in $\CN_{\g}$ have been completely classified (see \cite{CM} for example). It is known that every  closed $G$-orbit in $\g$ is $\check{G}$-stable (see \cite{MVW}). In \cite{Kato}, Kato considered an exotic nilpotent cone and derived the Deligne–Langlands theory for those exotic nilpotent orbits. To compute the local intersection cohomology of orbit closures in the exotic nilpotent cone, Achar and Henderson studied in \cite{AH08} the so-called ``enhanced nilpotent cone'' $\CN_{\gl_n(\C)}\times\C^{n\times 1}$ and classified its $\GL_n(\C)$-orbits.

  On the other hand, the  $G$-orbits in $\g\times E$ or $\CN_\g\times E$ have been studied in literature for various motivations. 
In \cite{RS}, for the purpose of proving the multiplicity one conjectures, Rallis and Schiffmann  studied the enhanced standard representation $\g\times E$ of $G$, and gave a criterion for a $G$-orbit in $\g\times E$ to be closed. In \cite{NO}, to generalize the results of \cite{Kato,AH08}, Nishiyama and Ohta determined regular semisimple $\GL_n(\C)$-orbits and the structure of the null cone in $\gl_n(\C)\times\left(\C^{n\times 1}\right)^p\times\left(\C^{1\times n}\right)^q$. In order to  generalize Ohta's conditions in \cite{Oht}, Nishiyama gave in \cite{Nis14} certain sufficient conditions for the map between orbit spaces induced by the inclusions of algebraic groups and varieties to be injective, and showed that the natural embedding of $\sp_{2n}(\C)\times\C^{2n\times 1}\hookrightarrow\gl_{2n}(\C)\times\C^{2n\times 1}\times\C^{1\times 2n}$ induces an injection of orbit spaces.

In what follows, we outline the main results of this paper. 
 
	\subsection{Closed orbits}
	\label{sec:closed}
	For $g_1\in\gl_{m_1}(\C),\dots,g_b\in\gl_{m_b}(\C)$, denote by $\diag(g_1,\dots,g_b)$ the matrix
	\begin{equation*}
		\begin{matrix}
			\left(\begin{matrix}
				g_1&&\\&\ddots&\\&&g_b
			\end{matrix}\right)\in\gl_{m_1+\dots+m_b}(\C).
		\end{matrix}
	\end{equation*}
	
	 To classify the closed $G$-orbits in $\g\times E$, we first study the closed $G$-orbits in $\CN_{\g}\times E$. For this purpose, we construct a subset $\fN_G$ of $\CN_{\g}\times E$ as follows:
    
   
	If $G=\GL_n(\C)$, then $\fN_G$ consists of all elements in $\CN_\g\times E$ which have the  form
		\begin{equation}
			\label{eq:closedorbitnullconeGL}
			x(k,y_1,\dots,y_k)=(\diag(J_k,0_{(n-k)\times(n-k)}),(y_1\dots,y_k,0\dots,0)^t,(1,0,\dots,0)),
		\end{equation}
where $k\geq 0$,  $y_1,\dots,y_k\in\C$ with $y_k\neq 0$, and \begin{equation*}
	J_0=0\quad\text{and}\quad	J_k=\left(\begin{matrix}
			0&1&&\\&\ddots&\ddots&\\&&0&1\\&&&0
		\end{matrix}\right)\ (k\ge 1).
	\end{equation*}
        
    If $G=\Sp_{2n}(\C)$, then $\fN_G$ consisting of all pairs $x=(X,u)\in\CN_\g\times E$ such that
		\begin{equation}
			\label{eq:nilpotentclosedspm}
			X=\left(\begin{matrix}
				\begin{array}{c:c}
					\diag(J_k,0_{(n-k)\times(n-k)})&\be_{kk}(n)\\\hdashline
					0&\diag(-J_k^t,0_{(n-k)\times(n-k)})
				\end{array}
			\end{matrix}\right)
		\end{equation}
		and
		\begin{equation}
			\label{eq:nilpotentclosedspv}
			u=(\underbrace{u_1,\dots,u_k}_k,\underbrace{0,\dots,0}_{n-k},\underbrace{u_{n+1},\dots,u_{n+k}}_k,\underbrace{0,\dots,0}_{n-k})^t,
		\end{equation}
        where $k\geq 0$, $u_1,\dots,u_k,u_{n+1},\dots,u_{n+k}\in\C$ with $u_{n+1}\neq 0$, and $\be_{i,j}(n)$ denotes the $n\times n$-matrix whose $(i,j)$-entry is 1 and other entries are zero. For $x\in\fN_G$ above, define $r(x)=k$.
        
        If $G=\RO_{2n+1}(\C)$, then $\fN_G$ consists of all pairs $(X,u)\in\CN_\g\times E$ such that
		\begin{equation}
			\label{eq:nilpotentclosedooddm}
			X=\left(\begin{matrix}
				\begin{array}{c:c:c:c:c}
					0&\begin{matrix}
						0&\dots&0
					\end{matrix}&\begin{matrix}
						0&\dots&0
					\end{matrix}&\begin{matrix}
						0&\ldots&0&1
					\end{matrix}&\begin{matrix}
						0&\dots&0
					\end{matrix}\\\hdashline
					\begin{matrix}
						0\\\vdots\\0\\-1
					\end{matrix}&J_k&0&0&0\\\hdashline
					\begin{matrix}
						0\\\vdots\\0
					\end{matrix}&0&0_{(n-k)\times(n-k)}&0&0\\\hdashline
					\begin{matrix}
						0\\\vdots\\0
					\end{matrix}&0&0&-J_k^t&0\\\hdashline
					\begin{matrix}
						0\\\vdots\\0
					\end{matrix}&0&0&0&0_{(n-k)\times(n-k)}
				\end{array}
			\end{matrix}\right)
		\end{equation}
		and
		\begin{equation}
			\label{eq:nilpotentclosedooddv}
		    u=(u_1,\underbrace{u_2,\dots,u_{k+1}}_k,\underbrace{0,\dots,0}_{n-k},\underbrace{u_{n+2},\dots,u_{n+k+1}}_k,\underbrace{0,\dots,0}_{n-k})^t,
		\end{equation}
        where $k\ge 0$, and $u_1,\dots,u_{k+1},u_{n+2},\dots,u_{n+k+1}\in\C$ with $u_{n+2}\neq 0$. For $x\in\fN_G$ above, define $r(x)=k$.
        
        If $G=\RO_{2n}(\C)$, then $\fN_G$ consists of all pairs $(X,u)\in\CN_\g\times E$ such that
		\begin{equation}
			\label{eq:nilpotentclosedoevenm}
			X=\left(\begin{matrix}
				\begin{array}{c:c}
					\diag(J_k,0_{(n-k)\times(n-k)})&\be_{k-1,k}(n)-\be_{k,k-1}(n)\\\hdashline
					0&\diag(-J_k^t,0_{(n-k)\times(n-k)})
				\end{array}
			\end{matrix}\right)
		\end{equation}
		and
		\begin{equation}
			\label{eq:nilpotentclosedoevenv}
			u=(\underbrace{u_1,\dots,u_k}_k,\underbrace{0,\dots,0}_{n-k},\underbrace{u_{n+1},\dots,u_{n+k}}_k,\underbrace{0,\dots,0}_{n-k})^t,
		\end{equation}
        where $k\ge 0$, and $u_1,\dots,u_k,u_{n+1},\dots,u_{n+k}$ with $u_{n+1}\neq 0$. For $x\in\fN_G$ above, define $r(x)=k$.

   

        The following theorem characterizes the closed $G$-orbits in $\CN_\g\times E$ whose proof is given in Section \ref{sec:4}.
	\begin{theorem}
		\label{thm:closedorbitsnullcone}
		For every $x\in\fN_G$, $Gx$ is a closed $G$-orbit in $\CN_\g\times E$. Conversely, each closed $G$-orbit in $\CN_{\g}\times E$ has such a form.
	\end{theorem}
    
For every $x,x'\in\fN_{G}$, we also give a sufficient and necessary condition for $Gx=Gx'$ (see Propositions \ref{prop:closedorbitglndifferent} and \ref{prop:closeddifferent}). For example, if $G=\GL_n(\C)$, then $Gx=Gx'$ if and only if $k=k'$ and $y_1=z_1,\dots,y_k=z_k$, where $x=x(k,y_1,\dots,y_k)$ and $x'=x(k',z_1,\dots,z_{k'})$ are as in \eqref{eq:closedorbitnullconeGL}.
    
Now we are going to describe the closed $G$-orbits in $\g\times E$. To this end, we construct a subset $\fX_G$ of $\g\times E$ as follows:

    If $G=\GL_n(\C)$, then $\fX_G$ consists of all elements in $\g\times E$ which have the form
    \begin{equation}
			\label{eq:closedorbitnilpotentgln}
			(\diag(c_1I_{n_1}+N_1,\dots,c_bI_{n_b}+N_b),(\left(u^{(1)}\right)^t,\dots,\left(u^{(b)}\right)^t)^t,(v^{(1)},\dots,v^{(b)})),
		\end{equation}
        where $n_1,\dots,n_b\geq 1$ such that $n_1+\dots+n_b=n$, $c_i\neq c_j$ for $1\leq i\neq j\leq b$, and $(N_i,u^{(i)},v^{(i)})\in\fN_{\GL_{n_i}(\C)}$.

    If $G=\Sp_{2n}(\C),\RO_{2n+1}(\C)$ or $\RO_{2n}(\C)$, then $\fX_G$ consists of all pairs $(X,u)\in\g\times E$ satisfying the following conditions:
    
    $\bullet$
    \begin{equation}
        \label{eq:closedsom1}
        X=\left(\begin{matrix}
			\begin{array}{c:c:c:c}
				N_0^{(1)}&&N_0^{(2)}&\\\hdashline
				&\diag(c_1I_{n_1}+N_1,\dots,c_bI_{n_b}+N_b)&&\\\hdashline
				&&N_0^{(3)}&\\\hdashline
				&&&\diag(-c_1I_{n_1}-N_1^t,\dots,-c_bI_{n_b}-N_b^t)
			\end{array}
		\end{matrix}\right),
     \end{equation}
            where $c_i\neq\pm c_j$ for $1\leq i\neq j\leq b$, $N_0^{(1)}\in\C^{n_0^{\prime}\times n_0^{\prime}}$, $N_0^{(2)}\in\C^{n_0^{\prime}\times n_0}$, $N_0^{(3)}\in\C^{n_0\times n_0}$, and $N_i\in\CN_{\gl_{n_i}(\C)}$ for $1\leq i\leq b$, with $n_0\geq 0$ and $n_0^{\prime},n_1,\dots,n_b\geq 1$ such that
		\begin{equation*}
			n_0+n_1+\dots+n_b=n\quad\text{and}\quad n_0^{\prime}=\begin{cases*}
				n_0+1,&if $G=\RO_{2n+1}(\C)$,\\
				n_0,&if $G=\Sp_{2n}(\C)$ or $\RO_{2n}(\C)$;
			\end{cases*}
		\end{equation*}
        
        $\bullet$
        \begin{equation}
        \label{eq:closedsov}
            u=\left(\left(u^{(0)}\right)^t,\left(u^{(1)}\right)^t,\dots,\left(u^{(b)}\right)^t,\left(v^{(0)}\right)^t,\left(v^{(1)}\right)^t,\dots,\left(v^{(b)}\right)^t\right)^t
        \end{equation}
        %with $u^{(0)}\in\C^{n_0^{\prime}\times 1},v^{(0)}\in\C^{n_0\times 1},u^{(1)},v^{(1)}\in\C^{n_1\times 1},\dots,u^{(b)},v^{(b)}\in\C^{n_b\times 1}$
        such that, for $1\leq i\leq b$, $(N_i,u^{(i)},\left(v^{(i)}\right)^t)\in\fN_{\GL_{n_i}(\C)}$, and
        \begin{equation*}
            \left(\left(\begin{matrix}
				N_0^{(1)}&N_0^{(2)}\\&N_0^{(3)}
			\end{matrix}\right),\left(\begin{matrix}
			    u^{(0)}\\v^{(0)}
			\end{matrix}\right)\right)\in\fN_{G_0},
        \end{equation*}
        where
        \begin{equation}
        	\label{eq:subgroup}
            G_0=\begin{cases*}
                \Sp_{2n_0}(\C),&if $G=\Sp_{2n}(\C)$,\\
                \RO_{2n_0+1}(\C),&if $G=\RO_{2n+1}(\C)$,\\
                \RO_{2n_0}(\C),&if $G=\RO_{2n}(\C)$.
            \end{cases*}
        \end{equation}


        Based on Theorem \ref{thm:closedorbitsnullcone}, we prove in Section \ref{sec:5} the following classification result. 
        
	\begin{theorem}
		\label{thm:closedorbit}
		For every $x\in\fX_G$, $Gx$ is a closed $G$-orbit in $\g\times E$. Conversely, each closed $G$-orbit in $\g\times E$ has such a form.
	\end{theorem}
    
In Section \ref{sec:5}, we also prove     
the following result, which says that the closed orbits in $\g\times E$ of  $\breve{G}$ coincide with that of $G$. 

    \begin{theorem}
		\label{thm:closedorbitsmvw}
        Every $G$-closed orbit in $\g\times E$ is $\breve{G}$-stable.
	\end{theorem}
	\subsection{Descents}
	%Applying Theorem~\ref{thm:closedorbit}, we determine all representations that may arise in the descendant of the enhanced standard representation.
    Let $O$ be a closed $\breve{G}$-orbit in $\g\times E$, and let $x\in O$. We denote by $\breve{G}_x$ the stabilizer of $\breve{G}$ at $x$, denote by $N_O^{\g\times E}$ the normal bundle of $O$ in $\g\times E$, and  denote by $N_{O,x}^{\g\times E}$ the fiber of $N_O^{\g\times E}$ at $x$.
    
	\begin{defn}
		\label{defn:descent}
		The natural action $\breve{G}_x$ on $N_{O,x}^{\g\times E}$ is called the \emph{descendant} of the enhanced standard representation at $x$.
	\end{defn}
%	\begin{theorem}
%		\label{thm:mainsimple}
%		Let $x\in\g\times E$ be such that $Gx$ is a closed orbit. The descendant $G_x\curvearrowright N_{Gx,x}^{\g\times E}$ must be as following:
%		
%		(1) Products of $\triv^{2k}$ and
%		\begin{equation*}
%			\GL_{k_1}(\C)\curvearrowright\gl_{k_1}(\C)\times\C^{k_1\times 1}\times\C^{1\times k_1},\dots,\GL_{k_b}(\C)\curvearrowright\gl_{k_b}(\C)\times\C^{k_b\times 1}\times\C^{1\times k_b}
%		\end{equation*}
%		with $k_1+\dots+k_b+k=n$, if $G=\GL_n(\C)$ and $E=\C^{n\times 1}\times\C^{1\times n}$;
%		
%		(2) Products of $\triv^{2k}$ and
%		\begin{align*}
%			&\Sp_{2l}(\C)\curvearrowright\sp_{2l}(\C)\times\C^{2l},\\
%			&\GL_{k_1}(\C)\curvearrowright\gl_{k_1}(\C)\times\C^{k_1\times 1}\times\C^{1\times k_1},\dots,\GL_{k_b}(\C)\curvearrowright\gl_{k_b}(\C)\times\C^{k_b\times 1}\times\C^{1\times k_b}
%		\end{align*}
%		with $l+k_1+\dots+k_b+k=n$, if $G=\Sp_{2n}(\C)$ and $E=\C^{2n}$;
%		
%		(3) Products of $\triv^k$ and
%		\begin{align*}
%			&\RO_{l}(\C)\curvearrowright\o_{l}(\C)\times\C^{l},\\
%			&\GL_{k_1}(\C)\curvearrowright\gl_{k_1}(\C)\times\C^{k_1\times 1}\times\C^{1\times k_1},\dots,\GL_{k_b}(\C)\curvearrowright\gl_{k_b}(\C)\times\C^{k_b\times 1}\times\C^{1\times k_b}
%		\end{align*}
%		with $l+2k_1+\dots+2k_b+k=m$, if $G=\RO_m(\C)$ and $E=\C^{m}$.
%	\end{theorem}

To describe such descendants, we need to define the MVW extension for a product of classical groups as well as its enhanced standard representation. 
Let $H_1,\dots,H_r$ be classical groups as in \eqref{eq:classicalgroups}, and set $H=H_1\times\dots\times H_r$. 
For $i=1,\dots,r$, write $\h_i\times E_i$ for the enhanced standard representation of $\breve{H}_i$ .

	\begin{defn}
		\label{defn:fiberproduct}
		We define the \emph{MVW extensions} $\breve{H}$ of $H$ to be the fiber product
		\begin{equation*}
			\breve{H}_1\times_{\{\pm 1\}}\dots\times_{\{\pm 1\}}\breve{H}_r:=\{(h_1,\dots,h_r,\delta)\mid (h_1,\delta)\in\breve{H}_1,\dots,(h_r,\delta)\in\breve{H}_r\}.
		\end{equation*}
       Additionally, we call the natural representation
        \begin{equation*}
            \h^{\mathrm{en}}=(\h_1\times E_1)\times\dots\times(\h_r\times E_r)
        \end{equation*}
         of $\breve{H}$ the  \emph{enhanced standard representation} of $\breve{H}$.
	\end{defn}
    
	%Let $\chi$ be the sign character from $\breve{H}$ to $\{\pm 1\}$ with kernel $H$. Namely, 
    %\[
    %\chi:\ \breve{H}\rightarrow \{\pm 1\},\quad (h_1,\dots,h_r,\delta)\mapsto \delta.
    %\]
 We denote by $\triv$ the trivial representation of $\breve{H}$. In the following result, we determine the descendants of enhanced standard representations $\g\times E$ of $\check{G}$. 
    
	\begin{theorem}
		\label{thm:mainmvwsimple}
		Let $O$ be a closed $\breve{G}$-orbit in $\g\times E$ and let $x\in O$. 
		
		\noindent(1) If $G=\GL_n(\C)$, then there exist $k\ge 0$ and $k_1,\dots,k_b\geq 1$ such that
		\begin{equation*}
		k+k_1+\dots+k_b=n,\quad	\breve{G}_x\simeq\breve{H},\quad \text{and}\quad 
            N_{O,x}^{\g\times E}\simeq \h^{\mathrm{en}}\oplus\triv^{2k},
		\end{equation*}
        where $H=\GL_{k_1}(\C)\times\dots\times\GL_{k_b}(\C)$.
		
		\noindent(2) If $G=\Sp_{2n}(\C)$, then there exist $l,k\geq 0$ and $k_1,\dots,k_b\geq 1$ such that
		\begin{equation*}
			k+l+k_1+\dots+k_b=n,\quad\breve{G}_x\simeq\breve{H}\quad \text{and}\quad 
            N_{O,x}^{\g\times E}\simeq \h^{\mathrm{en}}\oplus\triv^{2k},
		\end{equation*}
		where $H=\Sp_{2l}(\C)\times\GL_{k_1}(\C)\times\dots\times\GL_{k_b}(\C)$.
		
		\noindent(3) If $G=\RO_{n}(\C)$, then there exist $\gamma\in\{0,1\}$, $k,l\geq 0$ and $k_1,\dots,k_b\geq 1$ such that
		\begin{equation*}
			k+l+2k_1+\dots+2k_b+\gamma=n,\quad \breve{G}_x\simeq\breve{H}\quad \text{and}\quad 
            N_{O,x}^{\g\times E}\simeq \h^{\mathrm{en}}\oplus\triv^{2k+\gamma},
		\end{equation*}
		where $H=\RO_{l}(\C)\times\GL_{k_1}(\C)\times\dots\times\GL_{k_b}(\C)$.
	\end{theorem}
	
    The proof of Theorem \ref{thm:mainmvwsimple} is given in Section \ref{sec:6}.
    %We remark that the descendants of $\GL_n(\C)\curvearrowright\gl_n(\C)\times\C^{n\times 1}\times\C^{1\times n}$ appear naturally in the proof of the smooth transfer for Jacquet-Rallis trace formula. See \cite{Zha14, Xue19, CZ21} for details. Additionally, as pointed out in \cite{AGb}, by applying \cite{AG}*{Theorem 3.2.1},  Theorem \ref{thm:mainmvwsimple} can provide a direct proof of \cite{AGb}*{Proposition 3.2.1}, \cite{AGRS}*{Propositions 3.2, 5.2} and \cite{SZ}*{Propositions 7.1 and 7.2}.
	\section{Preliminaries}
    
\subsection{General notation}
\begin{itemize}
    \item  In this paper, all the (algebraic) varieties and 
groups are defined over $\C$.
\item We consider finite-dimensional vector spaces over $\C$ as algebraic varieties.
\item For the vector space $\C^{n\times 1}$, denote by $\{\be_1,\dots,\be_n\}$ its standard basis.
\item For an algebraic group $H$ acting on a variety $X$, a point $x\in X$, and a subset $K$ in $H$, we denote by
\begin{itemize}
\item  $X^H$ the set of all points in $X$ fixed by $H$,
\item $(H\bac X,\pi)$ the categorical quotient of $X$ by $H$ (if it exists), 
\item $H_x$ the stabilizer of $x$,
\item $Hx$ the $H$-orbit of $x$ in $X$, and
\item  $Z_H(K)$ the centralizer of $K$ in $H$.
\end{itemize}

\item For a Lie algebra $\h$ acting on a vector space $V$ and a vector $x\in V$, denote by $\h_x$ the stabilizer of $x$ in $\h$, and by $\h x$ the $\h$-orbit of $x$ in $V$.

\item For a variety $X$ and a point $x\in X$, denote by $T_xX$ the tangent space of $X$ at $x$. For a subvariety $Y$ of $X$ containing $x$, denote by $N_{Y,x}^X=(T_xX|_Y)/T_xY$ the normal space of $Y$ in $X$ at $x$.

\item For $y=(y_1,\dots,y_n)\in\C^{n\times 1}$, put
\begin{equation*}
	\phi(y)=\left(\begin{matrix}
		y_1&y_2&\dots&y_n\\
		&y_1&\dots&y_{n-1}\\
		&&\ddots&\vdots\\
		&&&y_1
	\end{matrix}\right)\quad\text{and}\quad\psi(y)=\left(\begin{matrix}
	    y_1&\dots&y_{n-1}&y_n\\
	    y_2&\dots&y_{n}&\\
	    \vdots&\iddots&&\\
	    y_n&&&
	\end{matrix}\right).
\end{equation*}
\end{itemize}
	\subsection{MVW extensions}
	Let $V$ be a vector space over $\C$. Denote by $\GL(V)$ the group of $\C$-linear automorphisms of $V$, and by $\gl(V)$ the Lie algebra of $\GL(V)$. The MVW-extension of $\GL(V)$ is
	\begin{equation*}
		\breve{\GL}(V)=\GL(V)\rtimes\{\pm 1\},
	\end{equation*}
	where $-1$ acts on $\GL(V)$ by transposition. 
      By specifying a basis of $V$, we have $\GL(V)=\GL_n(\C)$ and $\breve{\GL}(V)=\breve{\GL}_n(\C)$, where $n=\dim_{\C}V$.
	
	Assume now that $V$ is equipped with a quadratic or symplectic form $\langle\cdot,\cdot\rangle$. Put
	\begin{equation*}
		G(V)=\{g\in\GL(V)\mid\langle gv,gw\rangle=\langle v,w\rangle\text{ for }v,w\in V\}.
	\end{equation*}
	The Lie algebra of $G(V)$ is
	\begin{equation*}
		\g(V)=\{X\in\gl(V)\mid\langle Xv,w\rangle+\langle v,Xw\rangle=0\text{ for }v,w\in V\}.
	\end{equation*}
	Additionally, the MVW-extension $\breve{G}(V)$ of $G(V)$ is the subgroup of $\GL(V)\times\{\pm 1\}$ consisting of pairs $(g,\delta)$ such that either
	\begin{equation*}
		\delta=1\quad\text{and}\quad\langle gv,gw\rangle=\langle v,w\rangle\quad\text{for }v,w\in V,
	\end{equation*}
	or
	\begin{equation*}
		\delta=-1\quad\text{and}\quad\langle gv,gw\rangle=\langle w,v\rangle\quad\text{for }v,w\in V.
	\end{equation*}
    Let $\breve{G}(V)$ act on $\g(V)\times V$ by
    \begin{equation*}
        (g,\delta).(X,u)=(\delta gXg^{-1},\delta gu)
    \end{equation*}
    for $(g,\delta)\in\breve{G}(V)$, $X\in\g(V)$ and $u\in V$. This is a rational representation of $\breve{G}(V)$ (and $G(V)$ by taking restriction). We call this representation of $\breve{G}(V)$ (resp. $G(V)$) the \emph{enhanced standard representation} of $\breve{G}(V)$ (resp. $G(V)$).
    
   Note that if $G=\Sp_{2n}(\C)$ or $\RO_n(\C)$, then we have $G=G(E)$ and $\g=\g(E)$, where $E$ is equipped with the bilinear form $\langle\cdot,\cdot\rangle_E$ defined by
    \begin{equation*}
        \langle u,v\rangle_E=\begin{cases*}
            u^t\beta_{2n}v,&if $G=\Sp_{2n}(\C)$,\\
            u^t\alpha_nv,&if $G=\RO_n(\C)$.
        \end{cases*}
    \end{equation*}
    In this case, two definitions of the enhanced standard representation of $\breve{G}$ (resp. $G$) coincide.
	\subsection{Classical invariant theory}
	Let $H$ be a reductive group, acting on an affine variety $X$.
%	\begin{defn}
%		\label{defn:categoricalquotient}
%		The \emph{categorical quotient} of $X$ by $H$ is a pair $(Y,\pi)$, where $Y$ is analgebraic variety $Y$ and a morphism $\phi:X\to Y$ that is constant along $H$-orbits, such that for every other such pair $(Z,\phi)$, there exists a unique morphism $\phi^{\prime}:Y\to Z$ such that $\phi=\phi^{\prime}\circ\pi$.
%	\end{defn}
	It is known that the categorical quotient of $X$ by $H$ always exists. More precisely, $H\bac X=\Spec\left(\C[X]^H\right)$ and $\pi:X\to H\bac X$ is induced by the inclusion $\C[X]^H\hookrightarrow\C[X]$  (see \cite{PV}). Note that the morphism $\pi$ is surjective, and sends each $H$-invariant closed subset of $X$ onto a closed subset of $H\bac X$. Additionally,  every fiber of $\pi$ contains a unique closed orbit.
	
	\begin{theorem}[Luna's criterion, \cf \cite{PV}*{Remark of Theorem 6.17}]
		\label{thm:Lunacriterion}
		Let $K$ be a reductive subgroup of $H$, and let $x\in X^K$. Then the orbit $Hx$ is closed if and only if the orbit $Z_H(K)x$ is closed.
	\end{theorem}
    
	%The following result is useful to calculate $\dim H\bac X$.
%   Write $d=\max_{x\in X}\dim Hx$, and  set
%		\begin{equation*}
%			\Omega(X)=\{x\in X\mid Hx\text{ is closed and has dimension }d\}.
%		\end{equation*}
%	\begin{prop}
%		\label{cor:genericallyclosed}
%		Suppose that $X$ is irreducible and $\Omega(X)$ is nonempty. Then
%		\begin{equation*}
%			\dim H\bac X=\dim X-\dim H+\dim H_x
%		\end{equation*}
%		for all $x\in\Omega(X)$.
%	\end{prop}
%	\begin{proof}
%		Set $Y=\pi(\Omega(X))$. By \cite{FSR}*{Chapter 14, Theorem 3.13}, $\Omega(X)$ is an open subset in $X$, and so $Y$ is an open subset in $H\bac X$. Since $X$ is irreducible, $H\bac X$ is also irreducible, and then
%		\begin{equation*}
%			\dim X=\dim\Omega(X)\quad\text{and}\quad\dim H\bac X=\dim Y.
%		\end{equation*}
%		Note that the fibers of 
%		\begin{equation*}
%			\pi|_{\Omega(X)}:\Omega(X)\to Y
%		\end{equation*}
%		are closed orbits contained in $\Omega(X)$, and hence have the same dimension. Therefore
%		\begin{equation*}
%			\dim H-\dim H_x=\dim Hx=\dim\Omega(X)-\dim Y=\dim X-\dim H\bac X
%		\end{equation*}
%		for $x\in\Omega(X)$. This completes the proof.
%	\end{proof}
	\subsection{The algebra of invariants}
    
	The first step to classify closed orbits is to determine the algebra $\C[\g\times E]^{G}$ of invariants. 

    \begin{theorem}
		\label{thm:structure}
		The algebra $\C[\g\times E]^{G}$ is a polynomial ring with   \begin{equation*}
        \mathfrak A_G=\begin{cases*}
        \{\tr_1,\dots,\tr_n,\mu_0,\dots,\mu_{n-1}\},&if $G=\GL_{n}(\C)$,\\
            \{\tr_2,\tr_4,\dots,\tr_{2n},\eta_1,\eta_3,\dots,\eta_{2n-1}\},&if $G=\Sp_{2n}(\C)$,\\
            \{\tr_2,\tr_4,\dots,\tr_{2n},\eta_0,\eta_2,\dots,\eta_{2n}\},&if $G=\RO_{2n+1}(\C)$,\\
            \{\tr_2,\tr_4,\dots,\tr_{2n},\eta_0,\eta_2,\dots,\eta_{2n-2}\},&if $G=\RO_{2n}(\C)$.
        \end{cases*}
    \end{equation*} as a set of algebraic independent generators.
	\end{theorem}
	\begin{proof}
		When $G$ is a general linear group, this is \cite{NO}*{Theorem 2.1 (2)}. In the case of orthogonal or symplectic groups, the structure of $\C[\g\times E]^G$ is given in \cite{Sch}*{Table 3a.2, 3a.5 and Table 4a.3}.
	\end{proof}
Here, for every $i\ge 1$, 	$\tr_i$ denotes the polynomial on $\g\times E$ given by 
    \begin{equation*}
        \tr_i(X,u)=\tr(X^i).
    \end{equation*}
 For every $j\ge 0$, $\mu_j$ denotes the polynomial of $\gl_n(\C)\times \C^{n\times 1}\times \C^{1\times n}$ given by 
    \begin{equation*}
        \mu_j(X,u,v)=vX^ju.
    \end{equation*}
And, when $G=\Sp_{2n}(\C)$, $\RO_{2n+1}(\C)$ or $\RO_{2n}(\C)$, $\eta_j$ denotes the polynomial on $\g\times E$ defined by 
    \begin{equation*}
        \eta_j(X,u)=\langle X^ju,u\rangle_E.
    \end{equation*}
 
  
	\section{Closed orbits of the maximal dimension}
	
	In this section, we give a family of closed $G$-orbits in $\CN_\g\times E$, which has the maximal dimension.

   % The algebra $\C[\g\times E]^{G}$ is explicitly described in \cite{NO}*{Theorem 2.1} when $G$ is a general linear group, and in the case of orthogonal or symplectic groups, its structure is given in \cite{Sch}*{Table 3a.2, 3a.5 and Table 4a.3}.
   
    \subsection{The general linear case}
    %When $G=\GL_n(\C)$, Theorem \ref{thm:structure} is proved in \cite{NO}*{Theorem 2.1 (2)}. To prove the elements of $\mathfrak A_G$ are algebraically independent, Nishiyama and Ohta construct a family of closed $\GL_n(\C)$-orbits, which are  fibers of $\pi$. However, for the purpose of classifying the closed $\GL_n(\C)$-orbits in $\CN_{\gl_n(\C)}\times\C^{n\times 1}\times\C^{1\times n}$, we need to give another family of such closed orbits. Thus, we sketch a proof of Theorem \ref{thm:structure} (for the case that $G=\GL_n(\C)$) here.
    %\begin{prop}
    %    \label{prop:generators}
    %    The algebra $\C[\gl_n(\C)\times\C^{n\times 1}\times\C^{1\times n}]^{\GL_n(\C)}$ is generated by $$\tr_1,\dots,\tr_n,\mu_0,\dots,\mu_{n-1}.$$
    %\end{prop}
    %\begin{proof}
    %    This is proved in \cite{NO}*{Theorem 1.1}.
    %\end{proof}
    Assume that $G=\GL_n(\C)$. By Theorem \ref{thm:structure}, we regard the quotient $$\GL_n(\C)\bac(\gl_n(\C)\times\C^{n\times 1}\times\C^{1\times n})$$ as a closed subset of $\C^{1\times 2n}$, so that the quotient morphism is given by
	\begin{align*}
		\pi:\gl_n(\C)\times\C^{n\times 1}\times\C^{1\times n}&\to\GL_n(\C)\bac(\gl_n(\C)\times\C^{n\times 1}\times\C^{1\times n})\subseteq\C^{1\times 2n}\\
		x&\mapsto(\tr_1(x),\dots,\tr_n(x),\mu_0(x),\dots,\mu_{n-1}(x)).
	\end{align*}

    Recall the subset  $\fN_{\GL_n(\C)}$  of $\CN_{\gl_n(\C)}\times\C^{n\times 1}\times\C^{1\times n}$ defined in the Introduction.
    \begin{prop}
		\label{prop:maximalclosedgenerallinear}
		Let $x=x(n,y_1,\dots,y_n)\in\fN_{\GL_n(\C)}$. Then the stabilizer of $x$ is trivial, and $\GL_n(\C)x=\pi^{-1}(\underbrace{0,\dots,0}_n,y_1,\dots,y_n)$ is a closed orbit.
	\end{prop}
	\begin{proof}
		It is clear that $\pi(x)=(0,\dots,0,y_1,\dots,y_n)$ and that the stabilizer of $x$ is trivial. It remains to show that the fiber
		\begin{align*}
			O&=\pi^{-1}(0,\dots,0,y_1,\dots,y_n)\\
			&=\{(X,u,v)\in\gl_n(\C)\times\C^{n\times 1}\times\C^{1\times n}\mid X\text{ is nilpotent and }vu=y_1,\dots,vX^{n-1}u=y_n\}
		\end{align*}
		is an orbit.
		
		Let $(X,u,v)\in O$. Since $X$ is nilpotent and $X^{n-1}\neq 0$, we may assume that $X=J_n$. Write
		\begin{equation*}
			u=(u_1,\dots,u_n)^t\quad\text{and}\quad v=(v_1,\dots,v_n).
		\end{equation*}
		Then the condition that $vu=y_1,\dots,vJ_n^{n-1}u=y_n$ is equivalent to the equation
		\begin{equation*}
			\phi(v)u=(y_1,\dots,y_n)^t
		\end{equation*}
		In particular, $y_n=v_1u_n\neq 0$ forces that $v_1\neq 0$ and $u_n\neq 0$. Then $g=\phi(v)\in\GL_n(\C)$. It is easy to verify that
		\begin{equation*}
			gJ_ng^{-1}=J_n,\quad gu=(y_1,\dots,y_n)^t\quad\text{and}\quad(v_1,\dots,v_n)g^{-1}=(1,0,\dots,0).
		\end{equation*}
		Therefore the fiber $O$ is an orbit, as required.
	\end{proof}
	%\begin{cor}
	%	\label{cor:freegenerallinear}
	%	We have $\operatorname{Kdim}\C[\gl_n(\C)\times\C^{n\times 1}\times\C^{1\times n}]^{\GL_n(\C)}=2n$.
	%\end{cor}
	%\begin{proof}
	%	Let $x=x(n,y_1,\dots,y_n)\in\fN_{\GL_n(\C)}$. By Proposition \ref{prop:maximalclosedgenerallinear},  $\GL_n(\C)x$ is a closed orbit of maximal dimension. Then it follows from  Proposition \ref{cor:genericallyclosed} that
	%	\begin{align*}
	%		&\operatorname{Kdim}\C[\gl_n(\C)\times\C^{n\times 1}\times\C^{1\times n}]^{\GL_n(\C)}\\=&\dim(\gl_n(\C)\times\C^{n\times 1}\times\C^{1\times n})-\dim\GL_n(\C)=2n.
	%	\end{align*}
	%\end{proof}
%	\subsection{The generators for general linear Groups}
%	Denote by $\fS_s$ the symmetric group for $s\in\Z_{\geq 1}$. For any finite-dimensional vector space $V$, define the action of $\fS_s$ on $V^{\otimes s}$ by
%	\begin{equation*}
%		\sigma.(x_1\otimes\dots\otimes x_s)=x_{\sigma(1)}\otimes\dots\otimes x_{\sigma(s)},
%	\end{equation*}
%	for $x_1,\dots,x_s\in V$.
%	\begin{lm}[{\cite[\S 5.3.1]{GW}}]
%		\label{lm:multilinear}
%		Nonzero $\GL_n(\C)$-invariants of $\left(\C^{1\times n}\right)^{\otimes s}\otimes\left(\C^{n\times 1}\right)^{\otimes t}$ exist only for $s=t$. They are spanned by
%		\begin{align*}
%			\lambda_{\sigma}:\left(\C^{n\times 1}\right)^{\otimes s}\otimes\left(\C^{1\times n}\right)^{\otimes s}&\to\C,\\u_1\otimes\dots\otimes u_s\otimes v_1\otimes\dots\otimes v_s&\mapsto(v_1u_{\sigma(1)})\dots(v_su_{\sigma(s)})
%		\end{align*}
%		for $\sigma\in\fS_s$.
%	\end{lm}
%	Let $\tau=(i_1,\dots,i_k)\dots(j_1,\dots,j_m)\in\fS_r$ be the decomposition into disjoint cycles. We set
%	\begin{equation*}
%		\tr_{\tau}(X_1,\dots,X_r)=\tr(X_{i_1}\dots X_{i_k})\dots\tr(X_{j_1}\dots X_{j_m}).
%	\end{equation*}
%	\begin{lm}
%		\label{lm:multilineargl}
%		Nonzero $\GL_n(\C)$-invariants of $\left(\gl_n(\C)^*\right)^{\otimes r}\otimes\left(\C^{1\times n}\right)^{\otimes s}\otimes\left(\C^{n\times 1}\right)^{\otimes t}$ exist only for $s=t$. They are linearly generated by multi-linear functions $\gl_n(\C)^{\otimes r}\otimes\left(\C^{n\times 1}\right)^{\otimes s}\otimes\left(\C^{1\times n}\right)^{\otimes s}\to\C$ in the form of
%		\begin{equation}
%			\label{eq:multilinear}
%			\begin{aligned}
%				&\lambda:X_1\otimes\dots\otimes X_r\otimes u_1\otimes\dots\otimes u_s\otimes v_1\otimes\dots\otimes v_s\\
%				\mapsto&\tr_{\tau}(X_{i_1},\dots,X_{i_a})(v_{j_1}u_{j_{\sigma(1)}})\dots(v_{j_b}u_{j_{\sigma(b)}})\prod_c v_{k_c}(X_{m_1}\dots X_{m_d})u_{l_c}
%			\end{aligned}
%		\end{equation}
%		for $\tau\in\fS_{a}$ and $\sigma\in\fS_b$.
%	\end{lm}
%	\begin{proof}
%		There is an $\GL_n(\C)$-isomorphism
%		\begin{equation*}
%			\beta:\C^{n\times 1}\otimes\C^{1\times n}\to\gl_n(\C),\quad\beta(x\otimes y)u=(yu)x,\quad x,u\in\C^{n\times 1},y\in\C^{1\times n}.
%		\end{equation*}
%		Thus $\left(\gl_n(\C)^*\right)^{\otimes r}\otimes\left(\C^{1\times n}\right)^{\otimes s}\otimes\left(\C^{n\times 1}\right)^{\otimes t}\simeq\left(\C^{1\times n}\right)^{\otimes r+s}\otimes\left(\C^{n\times 1}\right)^{\otimes r+t}$. The first assertion follows from Lemma \ref{lm:multilinear}. The $\GL_n(\C)$-invariants of $\left(\C^{1\times n}\right)^{\otimes r+s}\otimes\left(\C^{n\times 1}\right)^{\otimes r+t}$ are spanned by multi-linear functions
%		\begin{align*}
%			&x_1\otimes\dots\otimes x_r\otimes u_1\otimes\dots\otimes u_s\otimes y_1\otimes\dots\otimes y_r\otimes v_1\otimes\dots\otimes v_s\\
%			\mapsto&(y_{i_1}x_{i_{\tau(1)}})\dots(y_{i_a}x_{i_{\tau(a)}})(v_{j_1}u_{j_{\sigma(1)}})\dots(v_{j_b}u_{j_{\sigma(b)}})\prod_c(v_{k_c}x_{m_1})(y_{m_1}x_{m_2})\dots(y_{m_d}u_{l_c}).
%		\end{align*}
%		Under the isomorphism $\beta$, these functions correspond to \eqref{eq:multilinear}.
%	\end{proof}
%	For $q_1,q_2,q_3\in\Z_{\geq 0}$ we say that $f\in\C[\gl_n(\C)\times\C^{n\times 1}\times\C^{1\times n}]$ is a homogeneous polynomial of degree $\bq=(q_1,q_2,q_3)$ if
%	\begin{equation*}
%		f(a_1X,a_2u,a_3v)=a_1^{q_1}a_2^{q_2}a_3^{q_3}f(X,u,v),\quad\text{for }a_1,a_2,a_3\in\C\smallsetminus\{0\}.
%	\end{equation*}
%	Denote by $\C[\gl_n(\C)\times\C^{n\times 1}\times\C^{1\times n}]_{\bq}$ the subspace of all homogeneous polynomials of degree $\bq$.
%	\begin{prop}
%		\label{prop:generators}
%		The algebra $\C[\gl_n(\C)\times\C^{n\times 1}\times\C^{1\times n}]^{\GL_n(\C)}$ is generated by
%		\begin{equation*}
%			\tr_1,\dots,\tr_n,\mu_0,\mu_1,\dots,\mu_{n-1}.
%		\end{equation*}
%	\end{prop}
%	\begin{proof}
%		Note that
%		\begin{align*}
%			\C[\gl_n(\C)\times\C^{n\times 1}\times\C^{1\times n}]_{\bq}&\simeq\Sym^{q_1}\left(\gl_n(\C)^*\right)\otimes\Sym^{q_2}(\C^{1\times n})\otimes\Sym^{q_2}(\C^{n\times 1})\\
%			&\simeq\left(\left(\gl_n(\C)^*\right)^{\otimes q_1}\otimes\left(\C^{1\times n}\right)^{\otimes q_2}\otimes\left(\C^{n\times 1}\right)^{\otimes q_3}\right)^{\fS_{q_1}\times\fS_{q_2}\times\fS_{q_3}}.
%		\end{align*}
%		Since the actions of $\GL_n(\C)$ and $\fS_{q_1}\times\fS_{q_2}\times\fS_{q_3}$ commute, we have
%		\begin{equation}
%			\label{eq:homogeneouspolygln}
%			\C[\gl_n(\C)\times\C^{n\times 1}\times\C^{1\times n}]_{\bq}^{\GL_n(\C)}\simeq\left(\left(\left(\gl_n(\C)^*\right)^{\otimes q_1}\otimes\left(\C^{1\times n}\right)^{\otimes q_2}\otimes\left(\C^{n\times 1}\right)^{\otimes q_3}\right)^{\GL_n(\C)}\right)^{\fS_{q_1}\times\fS_{q_2}\times\fS_{q_3}}.
%		\end{equation}
%		Applying Lemma \ref{lm:multilineargl}, we have $\C[\gl_n(\C)\times\C^{n\times 1}\times\C^{1\times n}]_{\bq}^{\GL_n(\C)}=0$ unless $q_2=q_3$. If $q_2=q_3$, the right side above is spanned by
%		\begin{equation*}
%			|\fS_{q_1}\times\fS_{q_2}\times\fS_{q_2}|^{-1}\sum_{(\tau,\sigma_1,\sigma_2)\in\fS_{q_1}\times\fS_{q_2}\times\fS_{q_2}}(\tau,\sigma_1,\sigma_2).\lambda,
%		where $\lambda$ are multilinear functions in the form \eqref{eq:multilinear}. Under the isomorphism \eqref{eq:homogeneouspolygln}, these tensors correspond to the polynomials
%		\begin{equation*}
%			F(X,u,v)=\lambda(X^{\otimes q_1}\otimes u^{\otimes q_2}\otimes v^{q_2})=\tr(X^{a_1})\dots\tr(X^{a_i})(vu)^{b}(vX^{d_1}u)\dots(vX^{d_j}u)
%		\end{equation*}
%		with $a_1+\dots+a_i+d_1+\dots+d_j=q_1$ and $b+j=q_2$. Since $\tr_i\in\C[\tr_1,\dots,\tr_n]$ for $i>n$ and $\mu_j\in\C[\mu_0,\mu_1,\dots,\mu_{n-1}]$ for $j>n-1$, $\C[\gl_n(\C)\times\C^{n\times 1}\times\C^{1\times n}]^{\GL_n(\C)}$ is generated by $\tr_1,\dots,\tr_n$ and $\mu_0,\mu_1,\dots,\mu_{n-1}$ as a $\C$-algebra.
%	\end{proof}
%	\subsection{The generators of $\C[\g\times E]^{G}$ for orthogonal and symplectic groups}
%In this and next subsections, we assume that $G=\Sp_{2n}(\C)$, $\RO_{2n+1}(\C)$ or $\RO_{2n}(\C)$. Note that in this case we have $G=G(E)$.

    
%	Let $k$ be a positive integer, and let $\rd=\{\{i_1,j_1\},\dots,\{i_k,j_k\}\}$ be a two-partition of the set $\{1,\dots,2k\}$. Define a $G$-invariant multilinear function $\lambda_{\rd}$ on $E^{\otimes 2k}$ by
%	\begin{equation*}
%		\lambda_{\rd}:E^{\otimes 2k}\to\C,\quad u_1\otimes\dots\otimes u_{2k}\mapsto\langle u_{i_1},u_{j_1}\rangle_E\dots\langle u_{i_k},u_{j_k}\rangle_E.
%	\end{equation*}
%	\begin{lm}[{\cite{GW}*{Theorem 5.3.5}}]
%		\label{lm:tensorinvariant}
%		Nonzero $G$-invariants of $\left(E^*\right)^{\otimes r}$ exist only if $r=2k$ is even, in which case they are spanned by $\lambda_{\rd}$, where $\rd$ runs over all two-partitions of $\{1,\dots,2k\}$.
%	\end{lm}
    
%	Consider the adjoint action of $G$ on $\gl(E)$. By the non-degenerated form $\langle\cdot,\cdot\rangle_E$, we  identify $E^*$ with $E$. Then there is a $G$-module isomorphism
%	\begin{equation*}
%		\beta:E\otimes E\xrightarrow{\sim}\gl(E).
%	\end{equation*}
%	such that $\beta(u\otimes v)x=\langle x,v\rangle_E u$ for $u,v,x\in E$. For $X\in\gl(E)$, denote by $X^*\in\gl(E)$ its adjoint operator, i.e., $\langle Xu,v\rangle_E=\langle u,X^*v\rangle_E$ for $u,v\in E$.

%   Fix positive integers $s,k$. Let $\rho:\{1,\dots,s\}\to\Z/2\Z$, and let $L=\{l_1,\dots,l_a\}\ (a\ge 1)$ be a subset of $\{1,\dots,s\}$.
%    Let $\tau\in\fS_a$, and pick a  cycle-decomposition $\tau=(i_1,\dots,i_p)\cdots(j_1,\dots,j_q)$ of it. We define
%    \begin{equation*}
%        \tr_{\rho,\tau}:\gl(E)^{\otimes a}\to\C,\quad(X_{l_1},\dots,X_{l_a})\mapsto\tr(Y_{l_{i_1}}\dots Y_{l_{i_p}})\dots\tr(Y_{l_{j_1}}\dots Y_{l_{j_q}}),
%    \end{equation*}
%    where
%    \begin{equation*}
%        Y_l=\begin{cases*}
%            X_l,&if $\rho(l)=1$,\\
%            X_l^*,&if $\rho(l)=-1$
%        \end{cases*}
%    \end{equation*}
%    for $l\in L$.

%    Let $\RP=\{L_1,\dots,L_k\}$ be a partition of $\{1,\dots,s\}\setminus L$, where we allow $L_i=\varnothing$. If $L_i=\{h_1^{(i)},\dots,h_{r_i}^{(i)}\}\neq\varnothing$, let $\sigma_i\in\fS_{r_i}$, and then define a multilinear map
%    \begin{equation*}
%        \gl(E)^{\otimes r_i}\to\gl(E),\quad(X_{h_1^{(i)}},\dots,X_{h_{r_i}^{(i)}})\mapsto Y_{\rho,i,\sigma_i}=Y_{h^{(i)}_{\sigma_i(1)}}\dots Y_{h^{(i)}_{\sigma_i(r_i)}},
%    \end{equation*}
%    where
%    \begin{equation*}
%        Y_h=\begin{cases*}
%            X_h,&if $\rho(h)=1$,\\
%            X_h^*,&if $\rho(h)=-1$
%        \end{cases*}
%    \end{equation*}
%    for $h\in L_i$. If $L_i=\varnothing$, we set $Y_{\rho,i,\sigma_i}=\id$.
    
%    Define $\Lambda$ to be the set consisting of multilinear functionals
%    \begin{align*}
%        \lambda_{\rho,L,\tau,\RP,\sigma,\rd}:\gl(E)^{\otimes s}\otimes E^{\otimes 2k}&\to\C,\\(X_1,\dots,X_s,u_1,\dots,u_{2k})&\mapsto\tr_{\rho,\tau}(X_{l_1},\dots,X_{l_a})\prod_{i=1}^k\langle Y_{\rho,i,\sigma_i}u_{e_i},u_{f_i}\rangle,
%    \end{align*}
%    where $\rho$, $L$, $\tau$ and $\RP$ are as above, $\sigma=(\sigma_1,\dots,\sigma_k)\in\fS_{r_1}\times\dots\times\fS_{r_k}$ and $\rd=\{\{e_1,f_1\},\dots,\{e_k,f_k\}\}$ is a two-partition of $\{1,\dots,2k\}$.
    
   
%	\begin{lm}
%		\label{lm:multilinearorthogonalsymplectic}
%		Nonzero $G$-invariants of $\left(\gl(E)^*\right)^{\otimes s}\otimes\left(E^*\right)^{\otimes t}$ exist only if $t=2k$ is even, in which case they are spanned by $\Lambda$.
%	\end{lm}
%	\begin{proof}
%		By the isomorphism $\beta$,  we see that $\left(\gl(E)^*\right)^{\otimes s}\otimes\left(E^*\right)^{\otimes t}$ is isomorphic to $(E^*)^{\otimes(2s+t)}$ as $G$-modules. Then the first assertion follows from Lemma \ref{lm:tensorinvariant}.
		
%		Assume that $t=2k$ is even. Then by Lemma \ref{lm:tensorinvariant}, $G$-invariants of $\left(E^*\right)^{\otimes(2s+2k)}$ are spanned by multilinear functionals
%		\begin{align*}
%			&x_1\otimes\dots\otimes x_s\otimes y_1\otimes\dots\otimes y_s\otimes u_1\otimes\dots\otimes u_{2k}\\
%			\mapsto&\langle x_{h_1},x_{h_2}\rangle_E\dots\langle x_{h_{2p-1}},x_{h_{2p}}\rangle_E\langle y_{h_1^{\prime}},y_{h_2^{\prime}}\rangle_E\dots\langle y_{h_{2p-1}^{\prime}},y_{h_{2p}^{\prime}}\rangle_E\langle x_{h_{p+1}},y_{h_{p+1}^{\prime}}\rangle_E\dots\langle x_{h_a},y_{h_a^{\prime}}\rangle_E\\
%			&\cdot\langle u_{i_1},u_{j_1}\rangle_E\dots\langle u_{i_b},u_{j_b}\rangle_E\prod_{c=b+1}^{k}\langle u_{i_c},z_{l_1}\rangle_E\langle z_{l_2},z_{l_3}\rangle_E\dots\langle z_{2l_q},u_{j_{c}}\rangle_E
%		\end{align*}
%		where $\{h_1,\dots,h_a\}=\{h_1^{\prime},\dots,h_a^{\prime}\}$ is a subset of $\{1,\dots,s\}$, $0\leq p\leq a/2$, $\{\{i_1,j_1\},\dots,\{i_k,j_k\}\}$ is a two-partition of $\{1,\dots,2k\}$, and, for $b+1\leq c\leq k$, $\{l_1,\dots,l_q\}$ are disjoint subsets of $\{1,\dots,s\}\setminus\{h_1,\dots,h_a\}$ and $(z_{l_1},\dots,z_{2l_q})$ is a permutation of $(x_{l_1},\dots,x_{l_q},y_{l_1},\dots,y_{l_q})$. Under the isomorphism $\beta$, these functionals correspond to functionals in $\Lambda$.
%	\end{proof}
%	For $q_1,q_2\in\Z_{\geq 0}$, we say that $f\in\C[\gl(E)\times E]^{G}$ is a homogeneous polynomial of degree $\bq=(q_1,q_2)$ if
%	\begin{equation*}
%		f(a_1X,a_2u)=a_1^{q_1}a_2^{q_2}f(X,u),\quad\text{for }a_1,a_2\in\C\setminus\{0\}.
%	\end{equation*}
%	Denote by $\C[\gl(E)\times E]_{\bq}$ the subspace of all homogeneous polynomials of degree $\bq$.

%    Denote by $\C\langle x,y\rangle$ the non-commutative polynomial ring in two variables. For $M\in\C\langle x,y\rangle$, define $\tr_M,\eta_M\in\C[\gl(E)\times E]^{G}$ by
%    \begin{equation*}
%        \tr_M(X,u)=\tr(M(X,X^*))\quad\text{and}\quad\eta_M(X,u)=\langle M(X,X^*)u,u\rangle.
%    \end{equation*}
%	\begin{prop}
%		\label{prop:generatororthogonalsymplectic}
%		The algebra $\C[\gl(E)\times E]^{G(E)}$ is generated by
%		\begin{equation*}
%			\{\tr_M,\eta_M\mid M\text{ is a monomial in }\C\langle x,y\rangle\}.
%		\end{equation*}
%	\end{prop}
%	\begin{proof}
%		First we have
%		\begin{equation*}
%			\C[\gl(E)\times E]_{\bq}\simeq\Sym^{q_1}\left(\gl(E)^*\right)\otimes\Sym^{q_2}(E^*)\simeq\left(\left(\gl(E)^*\right)^{\otimes q_1}\otimes(E^*)^{\otimes q_2}\right)^{\fS_{q_1}\times\fS_{q_2}}.
%		\end{equation*}
%		Since the action of $G$ and $\fS_{q_1}\times\fS_{q_2}$ commute, we have
%		\begin{equation}
%			\label{eq:homogegeouspolyorthgonalsymplectic}
%			\C[\gl(E)\times E]^{G}_{\bq}\simeq\left(\left(\left(\gl(E)^*\right)^{\otimes q_1}\otimes(E^*)^{\otimes q_2}\right)^{G}\right)^{\fS_{q_1}\times\fS_{q_2}}.
%		\end{equation}
%		By Lemma \ref{lm:multilinearorthogonalsymplectic}, $\C[\gl(E)\times E]_{\bq}^{G}=0$ unless $q_2$ is even. Assume that $q_2=2k$. Then the right side above is spanned by
%		\begin{equation*}
%			\left\{|\fS_{q_1}\times\fS_{q_2}|^{-1}\sum_{(\tau,\sigma)\in\fS_{q_1}\times\fS_{2k}}(\tau,\sigma).\lambda\mid\lambda\in\Lambda\right\}.
%		\end{equation*}
%		Under the isomorphism \eqref{eq:homogegeouspolyorthgonalsymplectic}, these functionals correspond to the polynomials
%		\begin{equation*}
%			F(X,u)=\lambda(X^{\otimes q_1}\otimes u^{\otimes 2k})=\prod_i\tr_{M_i}(X,u)\prod_j\eta_{M_j^{\prime}}(X,u),
%		\end{equation*}
%		where $M_i,M_j^{\prime}$ are monomials in $\C\langle x,y\rangle$. This completes the proof.
%	\end{proof}

   
%	\begin{cor}
%		\label{cor:generatororthogonalsymplectic}
%		The algebra $\C[\g(E)\times E]^{G}$ is generated by $\mathfrak A_{G}$.
%		
%		(1) $\tr_2,\tr_4,\dots,\tr_{2n},\eta_1,\eta_3,\dots,\eta_{2n-1}$, when $E=\C^{2n\times 1}$ is symplectic;
%		
%		(2) $\tr_2,\tr_4,\dots,\tr_{2n},\eta_2,\eta_4,\dots,\eta_{2n}$, when $E=\C^{(2n+1)\times 1}$ is orthogonal;
%		
%		(3) $\tr_2,\tr_4,\dots,\tr_{2n},\eta_2,\eta_4,\dots,\eta_{2n-2}$, when $E=\C^{2n\times 1}$ is orthogonal.
%	\end{cor}
%	\begin{proof}
%		Since $\g\times E$ is a $G$-invariant closed subset of $\gl(E)\times E$, the homomorphism
%		\begin{equation*}
%			\C[\gl(E)\times E]^{G}\to\C[\g\times E]^{G},
%		\end{equation*}
%		induced by the projection $\C[\gl(E)\times E]\to\C[\g\times E]$ is surjective. For $X\in\g$, we have $X^*=-X$. By Proposition \ref{prop:generatororthogonalsymplectic}, $\C[\g\times E]^{G}$ is generated by
%		\begin{equation*}
%			\tr_i,\ i\geq 1,\quad\text{and}\quad\eta_j,\ j\geq 1,
%		\end{equation*}
%		where $X\in\g$ and $u\in E$. Let $m=\dim_{\C}E$. Since $\tr_i\in\C[\tr_1,\dots,\tr_m]$ for $i>m$, and $\eta_j\in\C[\eta_0,\eta_1,\dots,\eta_{m-1}]$ for $j\geq m$, the algebra $\C[\g\times E]^{G}$ is generated by
%		\begin{equation*}
%			\tr_1,\dots,\tr_m,\eta_0,\eta_1,\dots,\eta_{m-1}.
%		\end{equation*}
		
%		When $E$ is symplectic, $\tr(X^i)=0$ if $i$ is odd, and $\langle X^ju,u\rangle_E=0$ if $j$ is even, for $(X,u)\in\g\times E$. When $E$ is orthogonal, $\tr(X^i)=0$ if $i$ is odd, and $\langle X^ju,u\rangle_E=0$ if $j$ is odd, for $(X,u)\in\g\times E$. Then we obtain the corollary.
%	\end{proof}
	\subsection{The symplectic and orthogonal cases}
    \label{section:closedmaximal}
	 Assume that $G=\Sp_{2n}(\C),\RO_{2n+1}(\C)$ or $\RO_{2n}(\C)$. Set $m=\dim_{\C}E$. Recall the map $r:\fN_G\to\Z_{\geq 0}$ defined in Subsection \ref{sec:closed}. Let $\fM_G$ be the set consisting of $x\in\fN_G$ such that $r(x)=n$.
	
	%If $G=\Sp_{2n}(\C)$, then $\fM_G$ consists of all pairs $(X,u)\in\fN_G$ such that
	%\begin{equation}
	%	\label{eq:symplecticclosedmaximal}
	%	X=\left(\begin{matrix}
	%		J_n&\be_{nn}(n)\\0&-J_n^t
	%	\end{matrix}\right)=\left(\begin{matrix}
	%	\begin{array}{c:c}
	%		\begin{matrix}
	%			0&1&&\\&\ddots&\ddots&\\&&0&1\\&&&0
	%		\end{matrix}&\begin{matrix}
	%			0&&&\\&\ddots&&\\&&0&\\&&&1
	%		\end{matrix}\\\hdashline
	%		0&\begin{matrix}
	%			0&&&\\-1&0&&\\&\ddots&\ddots&\\&&-1&0
	%		\end{matrix}
	%	\end{array}
	%	\end{matrix}\right)
	%\end{equation}
	%and
	%\begin{equation}
	%	\label{eq:symplecticclosedmaximalu}
	%	u=(u_1,\dots,u_{2n})^t\in\C^{2n\times 1},\quad\text{with }u_{n+1}\neq 0.
	%\end{equation}
	
	%If $G=\RO_{2n+1}(\C)$, then $\fM_G$ consists of all pairs $(X,u)\in\fN_G$ such that
	%\begin{equation}
	%	\label{eq:oddorthogonalclosedmaximal}
	%	X=\left(\begin{matrix}
	%		\begin{array}{c:c:c}
	%			0&0&\begin{matrix}
	%				0&\ldots&0&1
	%			\end{matrix}\\\hdashline
	%			\begin{matrix}
	%				0\\\vdots\\0\\-1
	%			\end{matrix}&J_n&0\\\hdashline
	%			0&0&-J_n^t
	%		\end{array}
	%	\end{matrix}\right)=\left(\begin{matrix}
	%		\begin{array}{c:c:c}
	%			0&0&\begin{matrix}
	%				0&\ldots&0&1
	%			\end{matrix}\\\hdashline
	%			\begin{matrix}
	%				0\\\vdots\\0\\-1
	%			\end{matrix}&\begin{matrix}
	%				0&1&&\\&\ddots&\ddots&\\&&0&1\\&&&0
	%			\end{matrix}&0\\\hdashline
	%			0&0&\begin{matrix}
	%				0&&&\\-1&0&&\\&\ddots&\ddots&\\&&-1&0
	%			\end{matrix}
	%		\end{array}
	%	\end{matrix}\right)
	%\end{equation}
	%and
	%\begin{equation}
	%	\label{eq:oddorthogonalclosedmaximalu}
	%	u=(u_1,\dots,u_{2n+1})^t\in\C^{(2n+1)\times 1},\quad\text{with }u_{n+2}\neq 0.
	%\end{equation}
	
	%If $G=\RO_{2n}(\C)$, then $\fM_G$ consists of all pairs $(X,u)\in\fN_G$ such that
	%\begin{equation}
	%	\label{eq:evenorthogonalclosedmaximal}
	%	X=\left(\begin{matrix}
	%		J_n&\be_{n-1,n}(n)-\be_{n,n-1}(n)\\&-J_n^t
	%	\end{matrix}\right)=\left(\begin{matrix}
	%		\begin{array}{c:c}
	%			\begin{matrix}
	%				0&1&&\\&\ddots&\ddots&\\&&0&1\\&&&0
	%			\end{matrix}&\begin{matrix}
	%				0&&&\\&\ddots&&\\&&0&1\\&&-1&0
	%			\end{matrix}\\\hdashline
	%			0&\begin{matrix}
	%				0&&&\\-1&0&&\\&\ddots&\ddots&\\&&-1&0
	%			\end{matrix}
	%		\end{array}
	%	\end{matrix}\right)
	%\end{equation}
	%and
	%\begin{equation}
	%	\label{eq:evenorthogonalclosedmaximalu}
	%	u=(u_1,\dots,u_{2n})^t\in\C^{2n\times 1}\quad\text{with }u_{n+1}\neq 0.
	%\end{equation}
	
	\begin{lm}
		\label{lm:stabilizermaximal}
		Let $x=(X,u)\in\fM_G$. If $G=\Sp_{2n}(\C)$ or $\RO_{2n+1}(\C)$, then $G_x$ is trivial. If $G=\RO_{2n}(\C)$, then $G_x=\{\pm 1\}$.
	\end{lm}
	\begin{proof}
		Let $V=\Span_{\C}\{X^iu\mid i\in\Z_{\geq 0}\}\subseteq E$. Then $\langle\cdot,\cdot\rangle_E|_V$ is non-degenerate and $G_x=G(V^{\bot})$. If $G=\Sp_{2n}(\C)$ or $\RO_{2n+1}(\C)$, then $V=E$ and hence $G_x$ is trivial. When $G=\RO_{2n}(\C)$, $V^{\bot}=\Span_{\C}\{\be_n-\be_{2n}\}$ is one-dimensional, and hence $G_x=\RO(V^{\bot})=\{\pm 1\}$.
	\end{proof}
	
	Set $d=m-n$. By Theorem \ref{thm:structure}, the quotient morphism $$\pi:\g\times E\to G\bac(\g\times E)=\C^{1\times m}$$ is given by
	\begin{equation*}
		\pi(x)=\begin{cases*}
			(\tr_2(x),\dots,\tr_{2n}(x),\eta_1(x),\dots,\eta_{2n-1}(x)),&if $m=2n$ and $G=\Sp_{2n}(\C)$,\\
			(\tr_2(x),\dots,\tr_{2n}(x),\eta_0(x),\dots,\eta_{2n}(x)),&if $m=2n+1$ and $G=\RO_{2n+1}(\C)$,\\
			(\tr_2(x),\dots,\tr_{2n}(x),\eta_0(x),\dots,\eta_{2n-2}(x)),&if $m=2n$ and $G=\RO_{2n}(\C)$.
		\end{cases*}
	\end{equation*}
	
	\begin{prop}
		\label{prop:closedmaximal}
		Let $y=(0,\dots,0,y_1,\dots,y_d)\in G\bac(\g\times E)$ with $y_d\neq 0$. Then $\pi^{-1}(y)$ is a closed $G$-orbit with a representative $x\in\fM_G$
	\end{prop}
	\begin{proof}
		We prove the proposition for $G=\Sp_{2n}(\C)$. The other cases are proved similarly. Let $(Y,v)\in\pi^{-1}(y)$ and write $v=(v_1,\dots,v_{2n})$. Then $Y$ is nilpotent. Since
		\begin{equation*}
			\langle Y^{2n-1}v,v\rangle_E=y_n\neq 0,
		\end{equation*}
		we have $Y^{2n-1}\neq 0$. By \cite{CM}*{Proposition 5.2.3 and recipe 5.2.2}, $Y$ lies in the nilpotent orbit corresponding to the partition $[2n]$, and then it is $\Sp_{2n}(\C)$-conjugate with some element $X$ which is as in \eqref{eq:nilpotentclosedspm} with $k=n$. Let $g\in\Sp_{2n}(\C)$ be such that $gYg^{-1}=X$ and let $u=gv$. Then
		\begin{equation*}
			(-1)^{n-1}u_{n+1}^2=\langle X^{2n-1}u,u\rangle_E=y_n\neq 0,
		\end{equation*}
		so $u_{n+1}\neq 0$. Now $$\dim\Sp_{2n}(\C).(Y,v)=\dim\Sp_{2n}(\C)x=\dim\Sp_{2n}(\C)=n(2n+1)$$ for every $(Y,v)\in\pi^{-1}(y)$, i.e., each orbit in $\pi^{-1}(y)$ has dimension $n(2n+1)$. So $\pi^{-1}(y)$ contains exactly one orbit, which implies the proposition.
	\end{proof}

	\begin{cor}
		\label{cor:closedmaximal}
		If $x\in\fM_G$, then $Gx$ is a closed $G$-orbit.
	\end{cor}
    
	\section{Proof of Theorem \ref{thm:closedorbitsnullcone}}\label{sec:4}
    In this section, we give a proof of Theorem \ref{thm:closedorbitsnullcone}.
    
	\subsection{Proof of Theorem \ref{thm:closedorbitsnullcone}: the general linear case}
	In this subsection, we assume that $G=\GL_n(\C)$. Recall the quotient morphism $$\pi:\gl_n(\C)\times\C^{n\times 1}\times\C^{1\times n}\to\GL_n(\C)\bac(\gl_n(\C)\times\C^{n\times 1}\times\C^{1\times n})=\C^{1\times 2n}.$$
	\begin{prop}
		\label{prop:closedorbitgln}
		Let $x=x(k,y_1,\dots,y_k)\in\fN_{\GL_n(\C)}$. Then $\GL_n(\C)x$ is the unique closed $\GL_n(\C)$-orbit in the fiber$$\pi^{-1}(\underbrace{0,\dots,0}_n,y_1,\dots,y_k,\underbrace{0,\dots,0}_{n-k}).$$
Furthermore, the stabilizer $\GL_n(\C)_x\simeq\GL_{n-k}(\C)$.
	\end{prop}
	\begin{proof}
		Let
		\begin{equation*}
			y=(\underbrace{0,\dots,0}_n,y_1,\dots,y_k,\underbrace{0,\dots,0}_{n-k}),\quad\text{with }y_k\neq 0.
		\end{equation*}
		Then $x\in\pi^{-1}(y)$. If $k=n$, the proposition follows from Proposition \ref{prop:maximalclosedgenerallinear}. If $k=0$, it is clear that $\{(0,0,0)\}$ is a closed orbit in $\pi^{-1}(0)$.
        
        Now assume that $1\leq k\leq n-1$. Let $H=\GL_n(\C)_x$. Then
		\begin{equation*}
			H=\left\{\diag(I_k,g)\mid g\in \GL_{n-k}(\C)\right\}\simeq\GL_{n-k}(\C).
		\end{equation*}
		We have $x\in(\gl_n(\C)\times\C^{n\times 1}\times\C^{1\times n})^H$ and
		\begin{equation*}
			Z_{\GL_n(\C)}(H)=\left\{\diag(g,cI_{n-k})\mid g\in\GL_k(\C)\text{ and }c\in\C^{\times}\right\}.
		\end{equation*}
		By Proposition \ref{prop:maximalclosedgenerallinear}, $Z_{\GL_n(\C)}(H)x$ is closed in the closed subset
		\begin{equation*}
			\left\{\left(\left(\begin{matrix}
				Y&\\&0_{(n-k)\times(n-k)}
			\end{matrix}\right),\left(\begin{matrix}
				u^{\prime}\\0_{(n-k)\times 1}
			\end{matrix}\right),(v^{\prime},0_{1\times(n-k)})\right)\mid Y\in\gl_k(\C),u^{\prime}\in\C^{k\times 1}\text{ and }v^{\prime}\in\C^{1\times k}\right\},
		\end{equation*}
		and hence $Z_{\GL_n(\C)}(H)x$ is closed. By Theorem \ref{thm:Lunacriterion}, $\GL_n(\C)x$ is a closed orbit.
	\end{proof}
	
    %Theorem \ref{thm:closedorbitsnullcone} for $G=\GL_n(\C)$ follows immediately from Proposition \ref{prop:closedorbitgln}.
    
    \begin{prop}
        \label{prop:closedorbitglndifferent}
        Let $x,x'\in\fN_{\GL_n(\C)}$. Then $\GL_n(\C)x=\GL_n(\C)x'$ if and only if $k=k'$ and $y_1=z_1,\dots,y_k=z_k$, where $x=x(k,y_1,\dots,y_k)$ and $x'=x(k',z_1,\dots,z_{k'})$.
    \end{prop}
    \begin{proof}
        This follows immediately from Theorem \ref{thm:structure}.
    \end{proof}
    
	\subsection{Proof of Theorem \ref{thm:closedorbitsnullcone}: the symplectic and orthogonal cases}
	In this subsection, we assume that $G=\Sp_{2n}(\C),\RO_{2n+1}(\C)$ or $\RO_{2n}(\C)$. Recall the quotient morphism
	\begin{equation*}
		\pi:\g\times E\to G\bac (\g\times E)=\C^{1\times m},
	\end{equation*}
	where $m=\dim_{\C}E$.
	
	In the case that $G=\RO_{2n}(\C)$, we need an easy lemma.
	
	\begin{lm}
		\label{lm:evenorthogonalrestriction}
		Assume that $G=\RO_{2n}(\C)$. Let $x=(X,u)\in\fM_G$ and let $V=\Span_{\C}\{X^iu\mid i\in\Z_{\geq 0}\}$. We have $X\in\g(V)$.
	\end{lm}
	\begin{proof}
		Since $X|_V\in\End_{\C}(V)$, we have $X|_{V^{\bot}}\in\End_{\C}(V^{\bot})$. It suffices to show that $X|_{V^{\bot}}=0$, which follows from $\langle Xv,v\rangle_E=0$ for $v\in V^{\bot}$ since $\dim_{\C}V^{\bot}=1$.
	\end{proof}
	
	\begin{prop}
		\label{prop:closedorbit}
		Let $x=(X,u)\in\fN_{G}$ with $r(x)=k$. Then $Gx$ is a closed orbit, and the stabilizer
		\begin{equation}
			\label{eq:stabilizer}
			G_x\simeq\begin{cases*}
				\Sp_{2(n-k)}(\C),&if $G=\Sp_{2n}(\C)$,\\
				\RO_{2(n-k)}(\C),&if $G=\RO_{2n+1}(\C)$,\\
				\RO_{2(n-k)+1}(\C),& if $G=\RO_{2n}(\C)$,
			\end{cases*}
		\end{equation}
		if $k\neq 0$.
	\end{prop}
	\begin{proof}
		If $k=0$, then $x=(0,0)$ and $Gx$ is obviously a closed orbit. If $k=n$, the proposition follows from Corollary \ref{cor:closedmaximal}.
		
		Assume that $1\leq k\leq n-1$. Let
		\begin{equation}
			\label{eq:subspace}
			W_k=\begin{cases*}
				\Span_{\C}\{\be_1,\dots,\be_k,\be_{n+1},\dots,\be_{n+k}\},&if $G=\Sp_{2n}(\C)$ or $\RO_{2n}(\C)$,\\
				\Span_{\C}\{\be_1,\be_2,\dots,\be_{k+1},\be_{n+1},\dots,\be_{n+k}\},&if $G=\RO_{2n+1}(\C)$,
			\end{cases*}
		\end{equation}
		and let $V_k=\Span_{\C}\{u,Xu,\dots,X^{m-1}u\}$. Then $\langle\cdot,\cdot\rangle_E$ is non-degenerated on $V_k$, and $E=V_k\oplus V_k^{\bot}$. Thus $G_x=G(V_k^{\bot})$. Note that $V_k=W_k$ if $G=\Sp_{2n}(\C)$ or $\RO_{2n+1}(\C)$, and $V_k$ is a subspace of $W_k$ of codimension 1 if $G=\RO_{2n}(\C)$. This implies \eqref{eq:stabilizer}.
		
		Write $H=G_x$. Then $x\in(\g\times E)^H$ and $Z_{G}(H)=G(V_k)$. By Lemma \ref{lm:evenorthogonalrestriction}, $x\in\g(V_k)\times V_k$. We have seen that $Z_{G}(H)x$ is a closed orbit in the closed subset $\g(V_k)\times V_k$ of $\g\times E$, so $Gx$ is a closed orbit by Theorem \ref{thm:Lunacriterion}.
	\end{proof}
	\begin{prop}
		\label{prop:closedorbit2}
		Every closed $G$-orbit in $\CN_{\g}\times E$ has a representative $x\in\fN_{G}$.
	\end{prop}
	\begin{proof}
		Let
		\begin{equation*}
			y=(\underbrace{0,\dots,0}_n,\underbrace{y_1,\dots,y_d}_d,0,\dots,0)\in\C^{1\times m}
		\end{equation*}
		with $d\geq 1$ and $y_d\neq 0$. Let
		\begin{equation*}
			k=\begin{cases*}
				d,&if $G=\Sp_{2n}(\C)$ or $\RO_{2n}(\C)$,\\
				d-1,&if $G=\RO_{2n+1}(\C)$.
			\end{cases*}
		\end{equation*}
		We need to show that $\pi^{-1}(y)$ contains an element $x=(X,u)\in\fN_G$ with $r(x)=k$. If $k=0$, then $y=0$ and $\{(0,0)\}$ is the closed orbit in $\pi^{-1}(0)$. If $k=n$, this follows from Proposition \ref{prop:closedmaximal}.
		
		Assume that $1\leq k\leq n-1$. Let $W_k$ be defined by \eqref{eq:subspace}. Then we have a natural embedding $\g(W_k)\hookrightarrow\g$. Denote by
		\begin{equation*}
			\pi_k:\g(W_k)\times W_k\to G(W_k)\bac(\g(W_k)\times W_k)=\C^{1\times (k+d)}
		\end{equation*}
		the quotient morphism. By Proposition \ref{prop:closedmaximal}, $\pi_k^{-1}(0,\dots,0,y_1,\dots,y_d)$ is a closed $G(W_k)$-orbit containing a representative $x'=(X^{\prime},u^{\prime})\in\fN_{G(W_k)}$ with such that $r(x')=k$. Write
		\begin{equation*}
			u^{\prime}=(u_1,\dots,u_d,u_{n+1},\dots,u_{n+k})^t.
		\end{equation*}
		Let $x=(X,u)\in\fN_G$ such that $r(x)=k$ and
		\begin{equation*}
			u=(u_1,\dots,u_d,\underbrace{0,\dots,0}_{n-k},u_{n+1},\dots,u_{n+k},\underbrace{0,\dots,0}_{n-k})^t.
		\end{equation*}
		Then $x=(X,u)\in\pi^{-1}(y)$.
	\end{proof}
    If $x\in\fN_{G}$ with $r(x)=k$, set
    \begin{equation}
    	\label{eq:invariant}
        \eta(x)=\begin{cases*}
        	(\eta_1(x),\eta_3(x),\dots,\eta_{2k-1}(x)),&if $G=\Sp_{2n}(\C)$,\\
        	(\eta_0(x),\eta_2(x),\dots,\eta_{2k}(x)),&if $G=\RO_{2n+1}(\C)$,\\
        	(\eta_0(x),\eta_2(x),\dots,\eta_{2k-2}(x)),&if $G=\RO_{2n}(\C)$.
        \end{cases*}
    \end{equation}
    \begin{prop}
        \label{prop:closeddifferent}
        Let $x,x'\in\fN_{G}$. Then $Gx=Gx'$ if and only if $r(x)=r(x')$ and $\eta(x)=\eta(x')$.
    \end{prop}
    \begin{proof}
        This follows from Theorem \ref{thm:structure}.
    \end{proof}
    
    We now finish the proof of Theorem \ref{thm:closedorbitsnullcone}.
    \begin{proof}[Proof of Theorem \ref{thm:closedorbitsnullcone}]
        The theorem follows from Propositions \ref{prop:closedorbitgln}, \ref{prop:closedorbit} and \ref{prop:closedorbit2}.
    \end{proof}
    
	\section{Proof of Theorems \ref{thm:closedorbit} and \ref{thm:closedorbitsmvw}}\label{sec:5}
This section is devoted to a proof of Theorems  \ref{thm:closedorbit} and \ref{thm:closedorbitsmvw}.	For every $X\in \g$, write  
	\begin{equation*}
		X=X_{\rs}+X_{\rn}
	\end{equation*}
    for the Jordan decomposition of $X$, where
	$X_{\rs}$ is semisimple and $X_{\rn}$ is nilpotent. 
Write $L_{X_{\rs}}=Z_{G}(X_{\rs})$ for the centralizer of $X_{\rs}$ in $G$, and write
	\begin{equation*}
		\CN_{X_{\rs}}=\{Y\in\CN_{\g}\mid [X_{\rs},Y]=0\}.
	\end{equation*}
    Note that $\CN_{X_{\rs}}$ is $L_{X_{\rs}}$-stable.
	\begin{lm}
		\label{lm:closedorbitnilpotent}
		Let $O$ be a closed orbit in $\g\times E$ and let $(X,u)\in O$. Set $O_{\rn}=L_{X_{\rs}}.(X_{\rn},u)$. Then $O_{\rn}$ is a closed $L_{X_{\rs}}$-orbit in $\CN_{X_{\rs}}\times E$, and
		\begin{equation}
			\label{eq:closedorbitnilpotent}
			O\cap \left((X_{\rs}+\CN_{X_{\rs}})\times E\right)=(X_{\rs},0)+O_{\rn}.
		\end{equation}
	\end{lm}
	\begin{proof}
	Note that	the closeness of $O_{\rn}$ follows from the equality \eqref{eq:closedorbitnilpotent}. Thus, we only need to prove  this equality. 
		
		It is obvious that $(X_{\rs},0)+O_{\rn}\subseteq O\cap \left((X_{\rs}+\CN_{X_{\rs}})\times E\right)$. On the other hand, let $Y\in\CN_{\g}$ be such that $(X_{\rs}+Y,v)\in O$. Then   $(X_{\rs}+Y,v)=g.(X_{\rs}+X_{\rn},u)$ for some $g\in G$. Note that
		\begin{equation*}
			gX_{\rs}g^{-1}+gX_{\rn}g^{-1}
		\end{equation*}
		is the Jordan decomposition of $X_{\rs}+Y$, which forces that $g\in L_{X_{\rs}}$ and $g.(X_{\rn},u)=(Y,v)$. Thus, $(X_{\rs}+Y,v)\in(X_{\rs},0)+O_{\rn}$ and so $O\cap \left((X_{\rs}+\CN_{X_{\rs}})\times E\right)\subseteq(X_{\rs},0)+O_{\rn}$, as required. 
	\end{proof}
	
	\subsection{Proof of Theorem \ref{thm:closedorbit}: the general linear case}
    In this subsection, we assume that $G=\GL_n(\C)$. First, by Lemma \ref{lm:closedorbitnilpotent}, we have the following result.
	\begin{prop}
		\label{prop:closedorbitnilpotentgln}
		Each closed $\GL_n(\C)$-orbit of $\gl_n(\C)\times\C^{n\times 1}\times\C^{1\times n}$ has a representative in $\fX_{\GL_n(\C)}$.
	\end{prop}
	%\begin{proof}
	%	Let $O$ be a closed $\GL_n(\C)$-orbit in $\gl_n(\C)\times\C^{n\times 1}\times\C^{1\times n}$, and let $(X,u,v)\in O$. We may assume that $X_{\rs}$ is the diagonal matrix without loss of generality. Then the proposition follows from Lemma \ref{lm:closedorbitnilpotent}.
	%\end{proof}
    
    In what follows, we are going to show that   $\GL_n(\C)x$ is  closed  for every $x\in\fX_{\GL_n(\C)}$.
    
	\begin{lm}
		\label{lm:regularsemisimplegln}
		Let $x=(X,u,v)\in\gl_n(\C)\times\C^{n\times 1}\times\C^{1\times n}$. If $\{u,Xu,\dots,X^{n-1}u\}$ is a basis of $\C^{n\times 1}$ and $\{v,vX,\dots,vX^{n-1}\}$ is a basis of $\C^{1\times n}$, then $\GL_n(\C)x$ is a closed orbit and the centralizer of $x$ is trivial.
	\end{lm}
	\begin{proof}
		This is proved in \cite{RS}*{Theorem 6.3}.
	\end{proof}
    
	\begin{prop}
		\label{prop:closedorbitmaximalgln}
		Let $x=(X,u,v)$ such that
		\begin{equation*}
			X=\diag(c_1I_{n_1}+J_{n_1},\dots,c_bI_{n_b}+J_{n_b})\in\gl_n(\C)
		\end{equation*}
		with $c_i\neq c_j$ for $1\leq i\neq j\leq b$,
		\begin{equation*}
			u=\left(\left(u^{(1)}\right)^t,\dots,\left(u^{(b)}\right)^t\right)^t\in\C^{n\times 1},\quad\text{and}\quad v=(v^{(1)},\dots,v^{(b)})\in\C^{1\times n}
		\end{equation*}
		with $u^{(i)}=(u_1^{(i)},\dots,u_{n_i}^{(i)})^t\in\C^{n_i\times 1}$ ($u_{n_i}^{(i)}\neq 0$) and $v^{(i)}=(1,0,\dots,0)\in\C^{1\times n_i}$ for $1\leq i\leq b$. Then $\GL_n(\C)x$ is a closed orbit.
	\end{prop}
	\begin{proof}
		Let $V=\Span_{\C}\{u,Xu,\dots,X^{n-1}u\}$ and $W=\Span_{\C}\{v,vX,\dots,vX^{n-1}\}$. We claim that $V=\C^{n\times 1}$ and $W=\C^{1\times n}$. The proposition follows from the claim by Lemma \ref{lm:regularsemisimplegln}.
		
		Now we prove the claim. Let
		\begin{align*}
			S_i&=\left\{(X-c_iI_n)^k\prod_{l=i+1}^b(X-c_lI_n)^{n_l}u\mid k=0,\dots n_i-1\right\},\\
			T_i&=\left\{\prod_{l=1}^{i-1}(X-c_lI_n)^{n_l}(X-c_iI_n)^ku\mid k=0,\dots,n_i-1\right\}
		\end{align*}
		for $1\leq i\leq b$. Put $S=\bigcup_{i=1}^bS_i$ and $T=\bigcup_{i=1}^bT_i$. Then $S\subseteq V$ and $T\subseteq W$. It is easy to verify that $S$ is a basis of $\C^{n\times 1}$ and $W$ is a basis of $\C^{1\times n}$. Therefore $V=\C^{n\times 1}$ and $W=\C^{1\times n}$.
		%
		%Consider the following vectors in $V$:
		%\begin{gather*}
		%	u=(\underbrace{*,\dots,*,u_{n_1}^{(1)}}_{n_1},\dots,\underbrace{*,\dots,*,u_{n_i}^{(i)}}_{n_i},\dots,\underbrace{*,\dots,*,u_{n_b}^{(b)}}_{n_b})^t,\\
		%	(X-c_bI_n)u=(\underbrace{*,\dots,*,(c_1-c_b)u_{n_1}^{(1)}}_{n_1},\dots,\underbrace{*,\dots,*,(c_i-c_b)u_{n_i}^{(i)}}_{n_i},\dots,\underbrace{*,\dots,*,u_{n_b}^{(b)}}_{n_b-1},0)^t,\\
		%	(X-c_bI_n)^2u=(\underbrace{*,\dots,*,(c_1-c_b)^2u_{n_1}^{(1)}}_{n_1},\dots,\underbrace{*,\dots,*,(c_i-c_b)^2u_{n_i}^{(i)}}_{n_i},\dots,\underbrace{*,\dots,*,u_{n_b}^{(b)}}_{n_b-2},0,0)^t,\\
		%	\vdots\\
		%	\begin{aligned}
		%		&\begin{aligned}
		%			(X-c_jI_n)^k&\prod_{l=j+1}^b(X-c_lI_n)^{n_l}u=(\underbrace{*,\dots,*,(c_1-c_j)^k\prod_{l=j+1}^b(c_1-c_l)^{n_l}u_{n_1}^{(1)}}_{n_1},\dots,\\&\underbrace{*,\dots,*,(c_i-c_j)^k\prod_{l=j+1}^b(c_i-c_l)^{n_l}u_{n_i}^{(i)}}_{n_i},\dots,\underbrace{*,\dots,*,\prod_{l=j+1}^b(c_j-c_l)^{n_l}u_{n_j}^{(j)}}_{n_j-k},\\
        %            &\underbrace{0,\dots,0}_{k},\underbrace{0,\dots,0}_{n_{j+1}+\dots+n_b})^t\quad\text{for }1\leq j\leq b-1\text{ and }0\leq k\leq n_j-1.
		%		\end{aligned}
		%	\end{aligned}
		%\end{gather*}
		%Since $u_{n_i}^{(i)}\neq 0$ and $c_i\neq c_j$ for $1\leq i\neq j\leq b$, these vectors form a basis of $\C^{n\times 1}$.
		%
		%Similarly, the vectors in $W$,
		%\begin{gather*}
		%	v,v(X-c_1I_n),v(X-c_1I_n)^2,\dots,v(X-c_1I_n)^{n_1},\\
		%	v\prod_{l=1}^{j-1}(X-c_lI_n)^{n_l}(X-c_jI_n),\dots,v\prod_{l=1}^{j-1}(X-c_lI_n)^{n_l}(X-c_jI_n)^{n_j}\quad(2\leq j\leq b-1),\\
		%	v\prod_{l=1}^{b-1}(X-c_lI_n)^{n_l}(X-c_bI_n),\dots,v\prod_{l=1}^{b-1}(X-c_lI_n)^{n_l}(X-c_bI_n)^{n_b-1},
		%\end{gather*}
		%form a basis of $\C^{1\times n}$.
	\end{proof}
	\begin{prop}
		\label{prop:closedgln}
		For every $x\in\fX_{\GL_n(\C)}$, the orbit $\GL_n(\C)x$ is closed.
	\end{prop}
	\begin{proof}
        Write $x=(X,u,v)$ as in \eqref{eq:closedorbitnilpotentgln}. Then for $1\leq i\leq b$, we have
        \begin{equation*}
			N_i=\diag(
				J_{k_i},0_{(n_i-k_i)\times(n_i-k_i)})\quad\text{for some }0\leq k_i\leq n_i,
        \end{equation*}
        $u^{(i)}=(u_1^{(i)},\dots,u_{k_i}^{(i)},0,\dots,0)^t\in\C^{n_i\times 1}$ ($u_{k_i}^{(i)}\neq 0$), and $v^{(i)}=(1,0,\dots,0)\in\C^{1\times n_i}$
    
		View $\C^{n\times 1}$ as the direct sum of $\C^{n_1\times 1},\dots,\C^{n_b\times 1}$. Denote by $\{\be_1^{(i)},\dots,\be_{n_i}^{(i)}\}$ the standard basis of $\C^{n_i\times 1}$. Then $\{\be_1^{(1)},\dots,\be_{n_1}^{(1)},\dots,\be_1^{(b)},\dots,\be_{n_b}^{(b)}\}$ is a standard basis of $\C^{n\times 1}$. Put
		\begin{equation*}
			H=\GL_n(\C)_x=\left\{\diag(I_{k_1},g_1,\dots,I_{k_b},g_b)\mid g_1\in\GL_{n_1-k_1}(\C),\dots,g_b\in\GL_{n_b-k_b}(\C)\right\}.
		\end{equation*}
		Then $x\in (\gl_n(\C)\times\C^{n\times 1}\times\C^{1\times n})^H$. Let $V=\Span_{\C}\{\be_1^{(1)},\dots,\be_{k_1}^{(1)},\dots,\be_1^{(b)},\dots,\be_{k_b}^{(b)}\}$, and $V_i=\Span_{\C}\{\be_{k_i+1}^{(i)},\dots,\be_{n_i}^{(i)}\}$ for $1\leq i\leq b$. Then
		\begin{equation*}
			Z_{\GL_n(\C)}(H)=\GL(V)\times\C^{\times}\cdot I_{V_1}\times\dots\times\C^{\times}\cdot I_{V_b}.
		\end{equation*}
		By Proposition \ref{prop:closedorbitmaximalgln}, $Z_{\GL_n(\C)}(H)x$ is a closed orbit in $\gl_n(\C)\times\C^{n\times 1}\times\C^{1\times n}$. Therefore $\GL_n(\C)x$ is a closed orbit by Theorem \ref{thm:Lunacriterion}.
	\end{proof}
	\subsection{Proof of Theorem \ref{thm:closedorbit}: the symplectic and orthogonal cases}
	Assume that $G=\Sp_{2n}(\C)$, $\RO_{2n+1}(\C)$ or $\RO_{2n}(\C)$. Set $m=\dim_{\C}E$. The following result is obvious.
	
	\begin{lm}
		\label{lm:0eigenspace}
		Let $X\in\g$. If the generalized 0-eigenspace $E_0$ of $X$ is nonzero, then the restriction of the bilinear form $\langle\cdot,\cdot\rangle_E$ on $E_0$ is non-degenerate. 
        %Furthermore, if $E$ is symplectic (resp. orthogonal), then $(E_0,\langle\cdot,\cdot\rangle|_{E_0})$ is symplectic (resp. o orthogonal).
	\end{lm}
%	\begin{proof}
%		Let $\pm c_1,\dots,\pm c_b$ be nonzero eigenvalues of $X$. Then
%		\begin{equation*}
%			E=E_0\oplus(E_1\oplus E_1^*)\oplus\dots\oplus(E_b\oplus E_b^*),
%		\end{equation*}
%		where $E_0$ is the 0-generalized eigenspace, and $E_i,E_i^*$ are $c_i,-c_i$-generalized eigenspace respectively. Then $E_0$ is orthogonal to $(E_1\oplus E_1^*)\oplus\dots\oplus(E_b\oplus E_b^*)$. Because $\langle\cdot,\cdot\rangle$ is non-degenerated, its restriction to $E_0$ has to be non-degenerated.
		
%		It is clear that if $E$ is symplectic (resp. orthogonal), then $(E_0,\langle\cdot,\cdot\rangle|_{E_0})$ is symplectic (resp. orthogonal). Since $\dim_{\C}E_i=\dim_{\C}E_i^*$, $\dim_{\C}E-\dim_{\C}E_0$ is even. Then we obtain the second assertion.
%	\end{proof}
By Lemma \ref{lm:closedorbitnilpotent}, we obtain the following result.

	\begin{prop}
		\label{prop:closedorbitnilpotent}
		Each closed $G$-orbit in $\g\times E$ has a representative $x\in\fX_G$.
	\end{prop}
	%\begin{proof}
	%	Let $O$ be a closed $G$-orbit in $\g\times E$ and let $(X,u)\in O$. We may assume that $X_{\rs}$ is a diagonal matrix in $\g$. Then the proposition follows from .
	%\end{proof}
    In what follows, we are going to show that   $Gx$ is  closed  for every $x\in\fX_{G}$.
    
    \begin{lm}
    	\label{lm:closedmaximal}
    	Let $x=(X,u)\in\g\times E$. If $\{u,Xu,\dots,X^{m-1}u\}$ is a basis of $E$, then $Gx$ is a closed orbit and the stabilizer of $x$ is trivial.
    \end{lm}
    \begin{proof}
    	This is proved in \cite{RS}*{Theorem 17.1}.
    \end{proof}
    
	\begin{lm}
		\label{lm:closedorbitmaximalso}
		Let $x=(X,u)\in\fX_G$. Write $X,u$ as in \eqref{eq:closedsom1} and \eqref{eq:closedsov}. Assume that $x$ satisfies the following conditions:
        \begin{itemize}
            \item $N_i=J_{n_i}$, $v^{(i)}=(1,0,\dots,0)$, $u^{(i)}=(u_1^{(i)},\dots,u_{n_i}^{(i)})^t$ with $u_{n_i}^{(i)}\neq 0$ for $1\leq i\leq b$,
            \item $\left(\left(\begin{matrix}
				N_0^{(1)}&N_0^{(2)}\\&N_0^{(3)}
			\end{matrix}\right),\left(\begin{matrix}
			    u^{(0)}\\v^{(0)}
			\end{matrix}\right)\right)\in\fM_{G_0}$, where $G_0$ is defined by \eqref{eq:subgroup}.
        \end{itemize}
		Let $V=\Span_{\C}\{u,Xu,\dots,X^{m-1}u\}$. Then $V$ is an orthogonal subspace of codimension 1, if $G=\RO_{2n}(\C)$ and $n_0\neq 0$. Otherwise $V=E$.
	\end{lm}
	\begin{proof}
		Let $E_0$ be the generalized 0-eigenspace of $X$, $E_i$ the generalized $c_i$-eigenspace, and $E_i^*$ the generalized $-c_i$-eigenspace, for $1\leq i\leq b$. Put
		\begin{equation*}
			E^{\prime}=(E_1\oplus E_1^*)\oplus\dots\oplus(E_b\oplus E_b^*).
		\end{equation*}
		Let $p$ be the projection from $E$ to $E'$ with respect to the decomposition $E=E_0\oplus E'$.
		
		Let $v=\prod_{i=1}^{b}(X+c_iI_n)^{n_i}u$,
		\begin{equation*}
			S_i^*=\left\{(X+c_iI_n)^k\prod_{l=i+1}^{b}(X+c_lI_n)^{n_l}u\mid k=0,\dots,n_i-1\right\},
		\end{equation*}
		and
		\begin{equation*}
			S_i=\left\{(X-c_iI_n)^k\prod_{l=i+1}^{b}(X-c_lI_n)^{n_l}v\mid k=0,\dots,n_i-1\right\}
		\end{equation*}
		for $1\leq i\leq b$. Put $S=\bigcup_{i=1}^b(S_i^*\cup S_i)$. By the proof of Proposition \ref{prop:closedgln}, $p(S)$ is a basis of $E^{\prime}$. Let
		\begin{equation*}
			u_0=\prod_{i=1}^b(X-c_iI_n)^{n_i}(X+c_iI_n)^{n_i}u
		\end{equation*}
		Then $u_0\in E_0$ and
		\begin{equation*}
			u_0=(s_1,\dots,s_{n_0^{\prime}},\underbrace{0,\dots,0}_{n_1+\dots+n_b},s_{n_0^{\prime}+1},\dots,s_{n_0^{\prime}+n_0},\underbrace{0,\dots,0}_{n_1+\dots+n_b})
		\end{equation*}
		with $s_{n_0^{\prime}+1}=(-1)^bc_1^2\dots c_b^2v_1^{(0)}\neq 0$.
		%
		%the images under $p$ of following vectors:
		%\begin{equation}
		%	\label{eq:closedmaximalso}
		%	\begin{gathered}
		%	    u,(X+c_bI_{n})u,\dots,(X+c_bI_{n})^{n_b}u,\dots,(X+c_1I_n)\prod_{i=2}^b(X+c_iI_n)^{n_i}u,\dots,\prod_{i=1}^b(X+c_iI_n)^{n_i}u,\\
		%	    (X-c_bI_n)\prod_{i=1}^b(X+c_iI_n)^{n_i}u,\dots,(X-c_bI_n)^{n_b}\prod_{i=1}^b(X+c_iI_n)^{n_i}u,\dots,\\
		%	    (X-c_1I_n)\prod_{i=2}^b(X-c_iI_n)^{n_i}\prod_{j=1}^b(X+c_jI_n)^{n_j}u,\dots,(X-c_1I_n)^{n_1-1}\prod_{i=2}^b(X-c_iI_n)^{n_i}\prod_{j=1}^b(X+c_jI_n)^{n_j}u
		%	\end{gathered}
		%\end{equation}
		
		Let $V=\Span_{\C}\{u,Xu,\dots,X^{m-1}u\}$ and $T=\{u_0,Xu_0,\dots,X^{n_0^{\prime}+n_0-1}u_0\}$. Assume that $G=\Sp_{2n}(\C)$ or $\RO_{2n+1}(\C)$. Then $T$ is a basis of $E_0$. Since $S$ and $T$ are contained in $V$, we have $V=E$.
		
		Assume that $G=\RO_{2n}(\C)$. If $n_0=0$, then $V=E$. Now assume that $n_0\geq 1$. Let $V_0=\Span_{\C}\{u_0,Xu_0,\dots,X^{2n_0-2}u_0\}$. Then $V=V_0\oplus E^{\prime}$, and hence $V^{\bot}$ equals the orthogonal complement of $V_0$ in $E_0$. Therefore $\dim_{\C}V^{\bot}=1$.
	\end{proof}
	\begin{prop}
		\label{prop:closedorbitso}
		For every $x\in\fX_G$, $Gx$ is a closed orbit.
	\end{prop}
	\begin{proof}
		Write $x=(X,u)$. Let $V=\Span_{\C}\{u,Xu,\dots,X^{m-1}u\}$. Then $x\in\g(V)\times V$. By Lemma \ref{lm:closedorbitmaximalso}, the restriction of the bilinear form $\langle\cdot,\cdot\rangle_E$ to $V$ is nondegenerate. Let $H=G_x$. Then $H=G(V^{\bot})$ and hence $Z_G(H)=G(V)$. By Lemma \ref{lm:closedmaximal}, $Z_G(H)x$ is a closed orbit in the closed subset $\g(V)\times V$. Therefore $Gx$ is a closed orbit by Theorem \ref{thm:Lunacriterion}.
	\end{proof}
    \subsection{Proof of Theorem \ref{thm:closedorbitsmvw}}
Here we give a proof of Theorem \ref{thm:closedorbitsmvw}, which states that every closed $G$-orbit in $\g\times E$ is $\breve{G}$-stable.

The following result is obvious.

	\begin{lm}
		\label{lm:mvwclosed}
		Let $H$ be a reductive group acting on an affine variety $X$, and let $K$ be a closed subgroup of $H$ which has index $2$. Then each closed $K$-orbit is $H$-stable if and only if $\C[X]^H=\C[X]^K$.
	\end{lm}
%	\begin{proof}
%		Let $\rho:K\bac X\to H\bac X$ be induced by the inclusion $\C[X]^H\subseteq\C[X]^K$. Then the following diagram commutes:
%		\begin{equation*}
%			\begin{tikzcd}
%				& X \arrow[ld, "\pi_{K\bac X}"'] \arrow[rd, "\pi_{H\bac X}"] &         \\
%				K\bac X \arrow[rr, "\rho"] &                           %                                & H\bac X
%			\end{tikzcd}
%		\end{equation*}
%		If each closed $K$-orbit is $H$-stable, $\rho$ is an isomorphism and hence $\C[X]^H=\C[X]^K$.
		
%		Assume that $\C[X]^H=\C[X]^K$. Then $\rho$ is the identity morphism. Let $Y=H\bac X=K\bac X$ and let $\pi=\pi_{H\bac X}=\pi_{K\bac X}$. Fix $y\in Y$. Let $x\in\pi^{-1}(y)$ such that $Kx$ is a closed orbit. Then
%		\begin{equation*}
%			Hx=\bigcup_{\dot{h}\in H/K}hHx
%		\end{equation*}
%		is a closed orbit in $\pi^{-1}(y)$. Since $[H:K]=2$, we have $hK=Kh$ for $h\in H$. Thus $hKx=Khx$ is a closed $K$-orbit in $\pi^{-1}(y)$ and hence $hKx=Kx$ for $h\in G$. Therefore $Hx=Kx$.
%	\end{proof}

In view of Lemma \ref{lm:mvwclosed}, Theorem \ref{thm:closedorbitsmvw} is implied by the following result. 
	\begin{lm}
		\label{lm:mvwinvariant}
		We have $\C[\g\times E]^{\breve{G}}=\C[\g\times E]^{G}$.
	\end{lm}
	\begin{proof}
		It follows immediately from Theorem \ref{thm:structure}.
		
        %By Theorem \ref{thm:structure}, we need to show that elements in $\mathfrak A_G$ are all $\breve{G}$-invariant. We now prove this case by case.
    
		%First assume that $G=\GL_n(\C)$. Then
		%\begin{equation*}
		%	\breve{G}=\breve{\GL}_n(\C)=\GL_n(\C)\rtimes\{\pm 1\}
		%\end{equation*}
        %Now
		%\begin{equation*}
		%	\left((I_n,-1).\tr_i\right)(X,u,v)=\tr(\left(X^t\right)^i)=\tr(X^i)=\tr_i(X,u,v)
		%\end{equation*}
		%and
		%\begin{equation*}
		%	\left((I_n,-1).\mu_j\right)(X,u,v)=(-u^t)\left(X^t\right)^j(-v^t)=(vX^ju)^t=vX^ju=\mu_j(X,u,v).
		%\end{equation*}
		%By Theorem \ref{thm:structure}, we have $\C[\gl_n(\C)\times\C^{n\times 1}\times\C^{1\times n}]^{\breve{\GL}_n(\C)}=\C[\gl_n(\C)\times\C^{n\times 1}\times\C^{1\times n}]^{\GL_n(\C)}$.
		
		%Assume that $G=\Sp_{2n}(\C)$. Then for $(g,-1)\in\breve{\Sp}_{2n}(\C)$, we have
		%\begin{equation*}
		%	\left((g,-1).\tr_{2i}\right)(X,u)=\tr(\left(-gXg^{-1}\right)^{2i})=\tr(gX^{2i}g^{-1})=\tr(X^{2i})=\tr_{2i}(X,u)
		%\end{equation*}
		%and
		%\begin{align*}
		%	\left((g,-1).\eta_{2j+1}\right)(X,u)&=\langle(-gXg^{-1})^{2j+1}(-gu),-gu\rangle=\langle-gX^{2j+1}u,gu\rangle\\
		%	&=-\langle u,X^{2j+1}u\rangle=\langle X^{2j+1}u,u\rangle=\eta_{2j+1}(X,u).
		%\end{align*}
		%Therefore $\C[\sp_{2n}(\C)\times\C^{2n\times 1}]^{\breve{\Sp}_{2n}(\C)}=\C[\sp_{2n}(\C)\times\C^{2n\times 1}]^{\Sp_{2n}(\C)}$ by Theorem \ref{thm:structure}.
		
		%Assume that $G=\RO_n(\C)$. Then we have
		%\begin{equation*}
		%	\left((I_n,-1).\tr_{2i}\right)(X,u)=\tr(\left(-X^{-1}\right)^{2i})=\tr(X^{2i})=\tr_{2i}(X,u)
		%\end{equation*}
		%and
		%\begin{equation*}
		%	\left((I_n,-1).\eta_{2j}\right)(X,u)=\langle(-X)^{2j}(-u),-u\rangle=\langle X^{2j}u,u\rangle=\eta_{2j}(X,u).
		%\end{equation*}
		%Therefore $\C[\o_{n}(\C)\times\C^{n\times 1}]^{\breve{\RO}_{n}(\C)}=\C[\o_{n}(\C)\times\C^{n\times 1}]^{\RO_{n}(\C)}$ by Theorem \ref{thm:structure}.
	\end{proof}
    
	
	
%    \begin{theorem}
%		\label{thm:mainsimple}
%		Let $x\in\g\times E$ be such that $Gx$ is a closed orbit. The descendant $G_x\curvearrowright N_{Gx,x}^{\g\times E}$ must be as following:
		
%		(1) If $G=\GL_n(\C)$, then there exist $k_1,\dots,k_b\geq 1$ with $k_1+\dots+k_b\leq n$ such that
%		\begin{equation*}
%			G_x\simeq H
%		\end{equation*}
 %       where $H=\GL_{k_1}(\C)\times\dots\times\GL_{k_b}(\C)$. Moreover, $N_{Gx,x}^{\g\times E}$ is isomorphic to the direct sum of $2k$ copies of the trivial representation, and the enhanced standard representation of $H$, where $k\geq 0$ such that $k_1+\dots+k_b+k=n$.
		
%		(2) If $G=\Sp_{2n}(\C)$, then there exist $l\geq 0$ and $k_1,\dots,k_b\geq 1$ with $l+k_1+\dots+k_b\leq n$ such that
%		\begin{equation*}
%			G_x\simeq H,
%		\end{equation*}
%		where $H=\Sp_{2l}(\C)\times\GL_{k_1}(\C)\times\dots\times\GL_{k_b}(\C)$. Moreover, $N_{Gx,x}^{\g\times E}$ is isomorphic to the direct sum of $2k$ copies of the trivial representation, and the enhanced standard representation of $\breve{H}$, where $k\geq 0$ such that $l+k_1+\dots+k_b+k=n$.
		
%		(3) If $G=\RO_{n}(\C)$, then there exist $l\geq 0$ and $k_1,\dots,k_b\geq 1$ with $2k_1+\dots+2k_b+l\leq n$ such that
%		\begin{equation*}
%			G_x\simeq H,
%		\end{equation*}
%		where $H=\RO_{l}(\C)\times\GL_{k_1}(\C)\times\dots\times\GL_{k_b}(\C)$. Moreover, $N_{Gx,x}^{\g\times E}$ is isomorphic to the direct sum of $2k+\gamma$ copies of the trivial representation, and the enhanced standard representation of $\breve{H}$,
%		where $\gamma\in\{0,1\}$ and $k\geq 0$ such that $l+2k_1+\dots+2k_b+2k+\gamma=n$.
%	\end{theorem}

	\section{Proof of Theorem \ref{thm:mainmvwsimple}}\label{sec:6}
    In this section, we give a proof of Theorem \ref{thm:mainmvwsimple}.

    Let $O$ be a closed $G$-orbit, and let $x\in O$. Let $\g_x$ be the Lie algebra of the stabilizer $G_x$. Then we have a $G_x$-module isomorphism
	\begin{equation*}
		T_xO=\g x\simeq\g/\g_x,
	\end{equation*}
	where $G_x$ acts by the adjoint action on $\g/\g_x$. Therefore the normal space
	\begin{equation}
		\label{eq:descent1}
		N_{O,x}^{\g\times E}\simeq\g_x\times E
	\end{equation}
	as rerpesentations of $G_x$, where $G_x$ acts on $\g_x$ by the adjoint action and acts on $E$ via the inclusion $G_x\subseteq G$.
	
	\subsection{Descents at $x\in\fN_{G}$ for $G=\GL_n(\C)$}
    In this subsection, we assume that $G=\GL_n(\C)$. Fix $x=x(k,y_1,\dots,y_k)\in\fN_{G}$.
    
    The following result is obvious.
    
    \begin{lm}
    	\label{lm:descentgln}
    	Let $h_1,h_2\in\GL_m(\C)$ be symmetric. For $i=1,2$, consider representations $(\rho_i,\C^{m\times 1}\times\C^{1\times m})$ and $(\omega_i,\gl_m(\C))$ of $\{\pm 1\}$ defined by
    	\begin{equation*}
    		\rho_i(-1)(u,v)=(-h_iv^t,-u^th_i^{-1})\quad\text{and}\quad\omega_i(-1)X=h_iX^th_i^{-1}.
    	\end{equation*}
    	Then
    	\begin{equation*}
    		\rho_1\simeq\rho_2,\quad (u,v)\mapsto(h_2h_1^{-1}u,v)
    	\end{equation*}
    	and
    	\begin{equation*}
    		\omega_1\simeq\omega_2,\quad X\mapsto h_2h_1^{-1}X.
    	\end{equation*}
    	In particular, $\rho_i\simeq\triv^m\oplus\sgn^m$ and $\omega_i\simeq\triv^{m(m+1)/2}\oplus\sgn^{(m-1)m/2}$ for $i=1,2$, where $\sgn$ denotes the sign character.
    \end{lm}
    
    Now we prove Theorem \ref{thm:mainmvwsimple} for $G=\GL_n(\C)$ and $x\in\fN_{G}$.
    \begin{prop}
    	\label{prop:descentglnmvw}
    	The stabilizer $\breve{\GL}_n(\C)_x\simeq\breve{\GL}_{n-k}(\C)$, and
    	\begin{equation}
    		\label{eq:descentgln}
    		N_{O,x}^{\g\times E}\simeq(\gl_{n-k}(\C)\times\C^{(n-k)\times 1}\times\C^{1\times(n-k)})\oplus\triv^{2k}
    	\end{equation}
    	as representations of $\breve{\GL}_{n-k}(\C)$.
    \end{prop}
    \begin{proof}
    	We have seen that
    	\begin{equation*}
    		\GL_n(\C)_x=\left\{\diag(I_k,g)\mid g\in\GL_{n-k}(\C)\right\}\simeq\GL_{n-k}(\C).
    	\end{equation*}
    	By Theorem \ref{thm:closedorbitsmvw}, we have $\breve{\GL}_n(\C)x=\GL_{n}(\C)x$, and hence $\GL_n(\C)_x$ is a subgroup of $\breve{\GL}_n(\C)_x$ of index two.
    	
    	Let $y=(y_1,\dots,y_k)$ and $g_0=\diag(-\psi(y),I_{n-k})\in\GL_n(\C)$. Then we have $(g_0,-1)\in\breve{\GL}_n(\C)_x$, and the map $\breve{\GL}_n(\C)_x\to\breve{\GL}_{n-k}(\C)$ defined by
    	\begin{equation*}
    		(\diag(I_k,g),1)\mapsto g\quad\text{and}\quad(g_0,-1)\mapsto(I_{n-k},-1).
    	\end{equation*}
    	is an isomorphism.
    	
    	Now we consider the action of $\breve{\GL}_{n-k}(\C)$ on $N_{O,x}^{\g\times E}$. For $z\in\g\times E$, write
    	\begin{equation*}
    		z=\left(\left(\begin{matrix}
    			A&B\\C&D
    		\end{matrix}\right),\left(\begin{matrix}
    		    u\\u'
    		\end{matrix}\right),(v,v')\right),
    	\end{equation*}
    	where $A\in\gl_k(\C),\ B\in\Mat_{k\times(n-k)}(\C),\ C\in\Mat_{(n-k)\times k}(\C),\ D\in\gl_{n-k}(\C),\ u\in\C^{k\times 1},\ u'\in\C^{(n-k)\times 1},\ v\in\C^{1\times k}$ and $v'\in\C^{1\times (n-k)}$. The action of $\breve{\GL}_{n-k}(\C)$ on $\g\times E$ is given by
    	\begin{equation*}
    		(g,1).z=\left(\left(\begin{matrix}
    			A&Bg^{-1}\\gC&gDg^{-1}
    		\end{matrix}\right),\left(\begin{matrix}
    			u\\gu'
    		\end{matrix}\right),(v,v'g^{-1})\right)
    	\end{equation*}
    	and
    	\begin{equation*}
    		(I_{n-k},-1).z=\left(\left(\begin{matrix}
    			\psi(y)A^t\psi(y)^{-1}&\psi(y)C^t\\B^t\psi(y)^{-1}&D
    		\end{matrix}\right),\left(\begin{matrix}
    			-\psi(y)v^t\\-v'^t
    		\end{matrix}\right),(-u^t\psi(y)^{-1},-u'^t)\right).
    	\end{equation*}
    	The tangent space
    	\begin{equation*}
    		T_xO=\left\{w=\left(\begin{matrix}
    			A&B\\C&0_{(n-k)\times(n-k)}
    		\end{matrix}\right)\mid A\in\gl_k(\C),\ B\in\Mat_{k\times(n-k)}(\C),\text{ and }C\in\Mat_{(n-k)\times k}(\C)\right\}
    	\end{equation*}
    	with the action of $\breve{\GL}_{n-k}(\C)$ given by
    	\begin{equation*}
    		(g,1).w=\left(\begin{matrix}
    			A&Bg^{-1}\\gC&0_{(n-k)\times(n-k)}
    		\end{matrix}\right)\quad\text{and}\quad(I_{n-k},-1).w=\left(\begin{matrix}
    		-\psi(y)A^t\psi(y)^{-1}&\psi(y)C^t\\B^t\psi(y)^{-1}&0_{(n-k)\times(n-k)}
    		\end{matrix}\right).
    	\end{equation*}
    	Comparing these two representations, we obtain \eqref{eq:descentgln} by Lemma \ref{lm:descentgln}.
    \end{proof}
    
    \subsection{Descents at $x\in\fN_{G}$ for $G=\Sp_{2n}(\C)$ or $\RO_m(\C)$}
    In this subsection, assume that $G=\Sp_{2n}(\C),\RO_{2n+1}(\C)$ or $\RO_{2n}(\C)$. Let $m=\dim_{\C}E$. Fix $x=(X,u)\in\fN_{G}$ with $r(x)=k$. Let $V=\Span_{\C}\{u,Xu,\dots,X^{m-1}u\}$ and let
    \begin{equation*}
    	(\gamma,\gamma')=\begin{cases*}
    		(0,1),&if $G=\Sp_{2n}(\C)$, or $(X,u)=(0,0)$,\\
    		(1,0),&if $G=\RO_{2n+1}(\C)$ and $(X,u)\neq(0,0)$,\\
    		(-1,0),&if $G=\RO_{2n}(\C)$ and $(X,u)\neq(0,0)$.
    	\end{cases*}
    \end{equation*}
    
    \begin{lm}
    	\label{lm:descent2}
    	Assume that $V\neq\{0\}$. Let $h'\in\g(V)$ be defined by
    	\begin{equation*}
    		h'(X^iu)=(-1)^{i+1}X^iu\quad\text{for }0\leq i\leq\dim_{\C}V-1.
    	\end{equation*}
    	Consider the representation of $\{\pm 1\}$ on $\g(V)$ given by $(-1).Y=-h'Yh'^{-1}$. Then
    	\begin{equation*}
    		\g(V)\simeq\triv^{k(k+\gamma+\gamma')}\oplus\sgn^{(k+\min\{\gamma,\gamma'\})^2}.
    	\end{equation*}
    \end{lm}
    \begin{proof}
    	We prove for $G=\Sp_{2n}(\C)$. The other cases are proved similarly. Let $J_e=\{0\leq i\leq 2k-1\mid i\text{ is even}\}$ and $J_o=\{0\leq i\leq 2k-1\mid i\text{ is odd}\}$. Identify $Y\in\g(V)$ with a matrix with respect to the basis $\{X^iu\mid i\in J_e\}\cup\{X^iu\mid i\in J_o\}$. Then
    	\begin{equation*}
    		\g(V)=\left\{\left(\begin{matrix}
    			Y_{1}&Y_{2}\\Y_{3}&Y_{4}
    		\end{matrix}\right)\mid Y_{i}\in\gl_k(\C),\ AY_1+Y_4^tA=AY_2-Y_2^tA=AY_3-Y_3^tA=0\right\},
    	\end{equation*}
    	where
    	\begin{equation*}
    		A=\left(\langle X^iu,X^ju\rangle\right)_{i\in J_e,j\in J_o}=\psi(\eta(x))\quad(\eta\text{ is defined by \eqref{eq:invariant}}),
    	\end{equation*}
    	and
    	\begin{equation*}
    		(-1).\left(\begin{matrix}
    			Y_{1}&Y_{2}\\Y_{3}&Y_{4}
    		\end{matrix}\right)=\left(\begin{matrix}
    		-Y_{1}&Y_{2}\\Y_{3}&-Y_{4}
    		\end{matrix}\right),
    	\end{equation*}
    	which implies the lemma.
    \end{proof}
    
    \begin{prop}
    	\label{prop:descentmvw}
    	The stabilizer $\breve{G}_x\simeq\breve{G}(V^{\bot})$ and 
    	\begin{equation}
    		\label{eq:descentmvw}
    		N_{O,x}^{\g\times E}\simeq(\g(V^{\bot})\times V^{\bot})\oplus\triv^{2k+\gamma},
    	\end{equation}
    	as representations of $\breve{G}(V^{\bot})$.
    \end{prop}
    \begin{proof}
    	If $x=(0,0)$ the proposition is trivial. Assume that $x\neq (0,0)$. Note that $\breve{G}_x$ preserves $V$ and $V^{\bot}$. We have a homomorphism
    	\begin{equation*}
    		\breve{G}_x\to\breve{G}(V^{\bot}),\quad (g,\delta)\mapsto (g|_{V^{\bot}},\delta),
    	\end{equation*}
    	which has an inverse
    	\begin{equation*}
    		\breve{G}(V^{\bot})\to\breve{G}_x\subseteq\GL(E)\times\{\pm 1\},\quad(h,\delta)\mapsto(\tilde{h},\delta)
    	\end{equation*}
    	such that
    	\begin{equation*}
    		\tilde{h}|_{V^{\bot}}=h\quad\text{and}\quad\tilde{h}(X^iu)=\delta^{i+1}X^iu,\text{ for }0\leq i\leq\dim_{\C}V-1.
    	\end{equation*}
    	Therefore $\breve{G}_x\simeq\breve{G}(V^{\bot})$.
    	
    	Now we consider the action of $\breve{G}(V^{\bot})$ on $N_{O,x}^{\g\times E}$. As vector spaces,
    	\begin{equation*}
    		\g\times E\simeq(\g(V)\times V)\oplus(\g(V^{\bot})\times V^{\bot})\oplus\Hom_{\C}(V,V^{\bot}),
    	\end{equation*}
    	and
    	\begin{equation*}
    		T_xO\simeq\g(V)\oplus\Hom_{\C}(V,V^{\bot}).
    	\end{equation*}
    	The actions of $\breve{G}(V^{\bot})$ are given by
    	\begin{equation}
    		\label{eq:descentmvw1}
    		(h,\delta).((Y,v),(Y',v'),Z)=((\delta h'Yh'^{-1},\delta h'v),(\delta hY'h^{-1},\delta hv'),Z)\quad\text{on }\g\times E
    	\end{equation}
    	and
    	\begin{equation}
    		\label{eq:descentmvw2}
    		(h,\delta).(Y,Z)=(h'Yh'^{-1},Z)\quad\text{on }T_xO,
    	\end{equation}
    	where $h'=\tilde{h}|_V$. Comparing \eqref{eq:descentmvw1} and \eqref{eq:descentmvw2}, we obtain \eqref{eq:descentmvw} by Lemma \ref{lm:descent2}.
    \end{proof}
    
    \subsection{Descents at general $x\in\fX_{G}$}
    Fix $x\in\fX_{G}$. Write $x$ in the form \eqref{eq:closedorbitnilpotentgln} if $G=\GL_n(\C)$, and write $x=(X,u)$ in the form \eqref{eq:closedsom1} and \eqref{eq:closedsov} if $G=\Sp_{2n}(\C),\RO_{2n+1}(\C)$ or $\RO_{2n}(\C)$. Set
	\begin{equation*}
		x_i=(N_i,u^{(i)},v^{(i)})\quad\text{for }1\leq i\leq b,
	\end{equation*}
	If $G=\Sp_{2n}(\C),\RO_{2n+1}(\C)$ or $\RO_{2n}(\C)$, let
	\begin{equation*}
		x_0=\left(\left(\begin{matrix}
			N_1^{(0)}&N_2^{(0)}\\&N_3^{(0)}
		\end{matrix}\right),\left(\begin{matrix}
			u^{(0)}\\v^{(0)}
		\end{matrix}\right)\right),
	\end{equation*}
	let $G_0$ be defined by \eqref{eq:subgroup} and then let $\g_0\times E_0$ be the enhanced standard representation of $G_0$. For convenience of notations, if $G=\GL_n(\C)$, let $x_0=0$, let $G_0$ be the trivial group, let $\breve G_0=\{\pm 1\}$, and let $\g_0\times E_0$ be the zero space.
	
	Now we have
	\begin{equation*}
		\breve{G}_x\simeq\left(\breve{G}_0\right)_{x_0}\times_{\{\pm 1\}}\breve{\GL}_{n_1}(\C)_{x_1}\times_{\{\pm 1\}}\dots\times_{\{\pm 1\}}\breve{\GL}_{n_b}(\C)_{x_b}.
	\end{equation*}
	It is clear that
	\begin{equation*}
		N_{G_0x_0,x_0}^{\g_0\times E_0}\times N_{\GL_{n_1}(\C)x_1,x_1}^{\gl_{n_1}(\C)\times\C^{n_1\times 1}\times\C^{1\times n_1}}\times\dots\times N_{\GL_{n_b}(\C)x_b,x_b}^{\gl_{n_b}(\C)\times\C^{n_b\times 1}\times\C^{1\times n_b}}
	\end{equation*}
	is a subrepresentation of $N_{O,x}^{\g\times E}$. By \eqref{eq:descent1} and comparing dimensions, we obtain
	\begin{equation}
		\label{eq:descent2}
		N_{O,x}^{\g\times E}\simeq N_{G_0x_0,x_0}^{\g_0\times E_0}\times N_{\GL_{n_1}(\C)x_1,x_1}^{\gl_{n_1}(\C)\times\C^{n_1\times 1}\times\C^{1\times n_1}}\times\dots\times N_{\GL_{n_b}(\C)x_b,x_b}^{\gl_{n_b}(\C)\times\C^{n_b\times 1}\times\C^{1\times n_b}}
	\end{equation}
	
    \begin{proof}[Proof of Theorem \ref{thm:mainmvwsimple}]
        The theorem follows from Propositions \ref{prop:descentglnmvw}, \ref{prop:descentmvw} and \eqref{eq:descent2}.
    \end{proof}
	\begin{thebibliography}{99}
		\bibitem[AGS]{AG}
		A. Aizenbud and D. Gourevitch, \textit{Generalized Harish-Chandra descent, Gelfand pairs, and an Archimedean analog of Jacquet-Rallis's theoremm, with an appendix by the authors and E. Sayag}, Duke Math. J., 149 (2009), no. 3, 509--567.
        \bibitem[AG]{AGb}
        A. Aizenbud and D. Gourevitch, \textit{Multiplicity one theorem for $(\GL_{n+1}(\R),\GL_n(\R))$}, Selecta Math., 15 (2009), 271--294.
        \bibitem[AGRS]{AGRS}
        A. Aizenbud, D. Gourevitch, S. Rallis and G. Schiffmann, \textit{Multiplicity one theorems}, Ann. Math. (2), 172 (2010), 1407--1434.  
        %\bibitem[AAG]{AAG}
        %A. Aizenbud, N. Avni and D. Gourevitch, \textit{Spherical pairs over close local fields}, Comment. Math. Helv. 87 (2012), no. 4, 929--962
		\bibitem[AH]{AH08}
		P. N. Achar and A. Henderson, \textit{Orbit closures in the enhanced nilpotent
		cone}, Adv. Math., 219 (2008), no. 1, 27--62. Correction in 228 (2011), no. 5 2984–2988.
		%\bibitem[AHJ]{AHJ}
		%P. N. Achar, A. Henderson, and B. F. Jones, \textit{Normality of orbit closures in the enhanced nilpotent cone}, Nagoya Math. J. 203 (2011), 1--45.
        \bibitem[BZSV]{BZSV}
         D. Ben-Zvi, Y. Sakellaridis and A. Venkatesh, \textit{Relative Langlands duality}, (2024), \eprint{arXiv:2409.04677}.
		\bibitem[CZ]{CZ21}
		P.-H. Chaudouard and M. Zydor, \textit{Le transfert singulier pour la formule des traces de Jacquet–Rallis}, Compos. Math. 157 (2021), no. 2, 303--434
		\bibitem[CM]{CM}
		D. Collingwood and W. McGovern, \textit{Nilpotent Orbits in Semisimple Lie Algebras}, Van Nostrand Reinhold Co., New York, 1993.
		\bibitem[FSR]{FSR}
		W. Ferrer Santos and A. Rittatore, \textit{Actions and Invariants of Algebraic Groups}, 2nd ed., Monographs and
		Research Notes in Mathematics, Chapman \& Hall/CRC, 2017.
        \bibitem[GGP]{GGP}
        W. T. Gan, B. H. Gross, and D. Prasad, \textit{Symplectic local root numbers, central critical L-values, and restriction problems in the representation theory of classical groups}, Ast\`erisque 346, (2012), 1--109.
        \bibitem[GK]{GK}
        I.M. Gelfand and D. Kazhdan, Representations of the group $\GL(n,K)$, where $K$ is a local field, Institute for Applied Mathematics, No. 942, 1971.
		\bibitem[GW]{GW}
		R. Goodment and N. Wallach, \textit{Symmetry, Representations, and Invariants}, Graduate Texts in Mathematics, 255. Springer-Verlag, Berlin-New York, 2009.
        \bibitem[JR]{JR}
        H. Jacquet and S. Rallis, \textit{On the Gross-Prasad conjecture for unitary groups}, in On Certain L-Functions, Clay Math. Proc. 13, Amer. Math. Soc., Providence, RI, 2011, 205–264.
		\bibitem[Kat]{Kato}
		S. Kato, \textit{An exotic Deligne–Langlands correspondence for symplectic groups}, Duke Math. J., 148 (2009), no. 2, 305--371.
		\bibitem[MVW]{MVW}
		C. Moeglin, M.-F. Vigneras, and J.-L. Waldspurger, Correspondence de Howe sur un corp p-adique, Lecture Notes in Mathematics, vol. 1291, Springer-Verlag, Berlin, 1987.
		\bibitem[Nis]{Nis14}
		K. Nishiyama, \textit{Enhanced orbit embedding}, Comment. Math. Univ. St. Pauli, 63 (2014), 223--232.
		\bibitem[NO]{NO}
		K. Nishiyama and T. Ohta, \textit{Enhanced adjoint actions and their orbits for the general linear group}, Pacific J. of Math., 298 (2019), no. 1, 141--155.
        \bibitem[Oht]{Oht}
        T. Ohta, \textit{An inclusion between sets of orbits and surjectivity of the restriction map of rings of invariants}, Hokkaido Math. J., 37 (2008), no. 3, 437–454.
		\bibitem[PV]{PV}
		V. L. Popov and E. B. Vinberg, \textit{Invariant Theory}, in \textit{Algebraic Geometry. IV}, A. N. Parshin and I. R. Shafarevich (eds.), Encyclopaedia of Mathematical Sciences, vol. 55, Springer-Verlag, Berlin, 1994, pp. 122--278.
		\bibitem[RS]{RS}
		S. Rallis and G. Schiffmann, \textit{Multiplicity one Conjectures}, \eprint{arXiv:0705.2168v1}.
		\bibitem[Sch]{Sch}
		G. W. Schwarz, \textit{Representations of simple Lie groups with regular rings of invariants}, Invent. Math., 49 (1978), 167--191.
        \bibitem[SZ1]{SZ1}
        B. Sun and C.-B. Zhu, \textit{A general form of Gelfand-Kazhdan criterion}, Manuscripta Math., 136 (2011), 185--197.
		\bibitem[SZ2]{SZ}
		B. Sun and C.-B. Zhu, \textit{Multiplicity on theorem: the Archimedean case}, Ann. of Math. (2), 175 (2012), 23--44.
		\bibitem[Sun]{Sun}
		B. Sun, \textit{Notes on MVW-extensions}, in \textit{Proceeding of the Fifth International Congress of Chinese Mathematicians}, L. Ji, Y. S. Poon, L. Yang and S.-T. Yau (eds.), AMS/IP Studies in Advanced Mathematics, vol. 51, part 1, American Mathematical Society, Providence, 2012, pp. 305--312.
		\bibitem[Xue]{Xue19}
		H. Xue, \textit{On the global Gan--Gross--Prasad conjecture for unitary groups: approximating smooth transfer of Jacquet--Rallis}, J. Reine Angew. Math. 756 (2019), 65--100.
		\bibitem[Zha]{Zha14}
		W. Zhang, \textit{Fourier transform and the global Gan-Gross-Prasad conjecture for unitary groups}, Ann. of Math. (2), 180 (2014), no. 3, 971--1049. 
	\end{thebibliography}
\end{document}
